% !TEX program = pdflatex
%=====================================================================
%  Criptografía y Seguridad --- ITBA · Finales resueltos
%  Documento maestro: preámbulo + \input de cada fragmento.
%=====================================================================
\documentclass[11pt,a4paper]{article}

\usepackage[utf8]{inputenc}
\usepackage[T1]{fontenc}
\usepackage[spanish,es-tabla,es-noshorthands]{babel}
\usepackage{lmodern}
\usepackage{microtype}

\usepackage{amsmath,amssymb,amsthm}
\usepackage{mathtools}
\usepackage{geometry}
\geometry{a4paper,margin=2.4cm,headheight=24pt}

\usepackage{xcolor}
%---- Paleta ITBA ----
\definecolor{itbaCyan}{HTML}{6EC6D9}
\definecolor{itbaCyanDark}{HTML}{2E8B9E}
\definecolor{itbaBlue}{HTML}{AEDCEA}
\definecolor{boxBlue}{HTML}{BBD4F0}
\definecolor{slideGray}{HTML}{ECECEC}
\definecolor{practGreen}{HTML}{4F8A3D}
\definecolor{practGreenL}{HTML}{E2EFD9}
\definecolor{attackRed}{HTML}{C0392B}
\definecolor{codeBg}{HTML}{F5F5F5}

\usepackage{graphicx}
\usepackage{enumitem}
\usepackage{booktabs}
\usepackage{array}
\usepackage{longtable}
\usepackage{tcolorbox}
\tcbuselibrary{skins,breakable}
\usepackage{fancyhdr}
\usepackage{titlesec}
\usepackage{listings}
\usepackage{caption}
\usepackage{pifont}
\usepackage{tikz}
\usetikzlibrary{shapes.geometric,shapes.misc,shapes.symbols,arrows.meta,positioning,
                calc,fit,backgrounds,decorations.pathmorphing,decorations.pathreplacing,
                decorations.markings,shadows.blur,chains,matrix}
\usepackage{hyperref}
\hypersetup{colorlinks=true,linkcolor=itbaCyanDark,urlcolor=itbaCyanDark,
            citecolor=itbaCyanDark,linktoc=all,
            pdftitle={Finales resueltos - Criptografia y Seguridad},
            pdfauthor={Criptografia y Seguridad - ITBA}}

%---- Encabezados ----
\pagestyle{fancy}
\fancyhf{}
\lhead{\small\textsc{Criptografía y Seguridad --- ITBA}}
\rhead{\small Finales resueltos}
\cfoot{\small\thepage}
\renewcommand{\headrulewidth}{0.4pt}

%---- Títulos de sección con barra de color ----
\titleformat{\section}
  {\normalfont\Large\bfseries\color{itbaCyanDark}}
  {\thesection}{0.6em}{}
  [\vspace{-0.4em}{\color{itbaCyan}\titlerule[2pt]}]
\titleformat{\subsection}
  {\normalfont\large\bfseries\color{itbaCyanDark}}
  {\thesubsection}{0.6em}{}

%---- Estilo de código ----
\lstdefinestyle{code}{
  basicstyle=\ttfamily\footnotesize, backgroundcolor=\color{codeBg},
  keywordstyle=\color{itbaCyanDark}\bfseries, commentstyle=\color{practGreen}\itshape,
  stringstyle=\color{attackRed}, breaklines=true, frame=single,
  rulecolor=\color{black!30}, showstringspaces=false, columns=fullflexible, tabsize=2,
  xleftmargin=10pt, framexleftmargin=8pt}
\lstset{style=code}

%=====================================================================
%  Entorno requerido por los fragmentos: consigna textual del examen
%=====================================================================
\newtcolorbox{consigna}{
  breakable, enhanced, colback=slideGray, colframe=black!55, boxrule=0.6pt, arc=2pt,
  left=8pt,right=8pt,top=6pt,bottom=6pt, fonttitle=\bfseries, coltitle=white,
  title={Consigna (textual)}}

% Cajas auxiliares
\newtcolorbox{defbox}[1]{enhanced,breakable,colback=itbaBlue!25,colframe=itbaCyanDark,
  boxrule=1pt,arc=2pt,fonttitle=\bfseries,coltitle=white,title={#1}}
\newtcolorbox{warnbox}[1]{enhanced,breakable,colback=attackRed!8,colframe=attackRed,
  boxrule=1pt,arc=2pt,fonttitle=\bfseries,coltitle=white,title={#1}}
\newtcolorbox{practbox}[1]{enhanced,breakable,colback=practGreenL,colframe=practGreen,
  boxrule=1pt,arc=2pt,fonttitle=\bfseries,coltitle=white,title={#1}}
\newtcolorbox{ojo}[1][]{breakable,enhanced,colback=orange!10!white,colframe=orange!80!black,
  fonttitle=\bfseries,title={\faWarning Cuidado --- error típico},coltitle=black,#1}
\newcommand{\faWarning}{$\triangle$\ }

%---- Macros de control de acceso ----
\newcommand{\dom}{\ensuremath{\succeq}}        % relación de dominancia
\newcommand{\Lvl}{\ensuremath{\mathrm{L}}}     % nivel de seguridad L(·)
\newcommand{\il}{\ensuremath{\mathrm{i}}}      % nivel de integridad i(·)
\providecommand{\cmark}{\ding{51}}
\providecommand{\xmark}{\ding{55}}

% Cita textual del profesor (tomada del sistema de resúmenes por unidad)
\newcommand{\prof}[1]{\medskip\noindent\textit{``#1''}\\[-2pt]\hspace*{1em}{\footnotesize\textcolor{gray}{--- dicho en clase}}\medskip}

%=====================================================================
\begin{document}
%=====================================================================

\begin{titlepage}
\centering
\vspace*{2cm}
{\LARGE\bfseries\color{itbaCyanDark} Finales resueltos de\\[3pt]
Criptografía y Seguridad}\\[0.5cm]
{\large Exámenes finales (72.44) --- resoluciones completas}\\[0.9cm]

{\normalsize
\begin{minipage}{0.88\textwidth}\itshape
Este PDF reúne la resolución de varios exámenes finales de la materia. Cada ejercicio
incluye la \textbf{consigna textual}, los \textbf{temas} involucrados (con referencia a la
clase/unidad de los resúmenes de la materia), \textbf{qué necesitás saber} y la
\textbf{resolución completa y razonada}, construida exclusivamente a partir del material
de la materia (resúmenes de Primera y Segunda Parte).
\end{minipage}}\\[1.0cm]

{\normalsize Instituto Tecnológico de Buenos Aires (ITBA)}
\vfill
{\footnotesize\itshape Generado a partir de los resúmenes de la materia y los enunciados de los finales.}
\end{titlepage}

\tableofcontents
\newpage

%=====================================================================
\section*{Introducción}
\addcontentsline{toc}{section}{Introducción}
Este documento resuelve los ejercicios de los siguientes exámenes finales de
Criptografía y Seguridad (72.44): \textbf{21/07/2021}, \textbf{17/07/2023},
\textbf{06/07/2023}, \textbf{10/07/2023}, \textbf{2025 --- C2 (1)}, y un conjunto de
\textbf{ejercicios adicionales} (bingo virtual, sensores de robo de cable, app
``Organizame''). Por cada ejercicio se muestra la consigna textual, los temas que toca
citando la unidad/clase correspondiente de los resúmenes de la materia (Primera Parte:
criptografía; Segunda Parte: seguridad), y una resolución completa.

\medskip
\noindent\textbf{Regla de oro seguida en toda resolución:} no se agrega nada que no esté
en el material de la materia (resúmenes de Primera y Segunda Parte). Cuando algún punto de
un enunciado excede ese material, se lo señala explícitamente en la resolución en lugar de
inventar contenido.

\bigskip
\noindent\textbf{Nota sobre las fuentes de los finales:} varios de los PDF de los exámenes
ya traían respuestas manuscritas de un alumno rendido. Esas respuestas se usaron solo como
referencia de qué se preguntaba; cada resolución de este documento fue re-elaborada y
verificada contra los resúmenes de la materia, seleccionando lo correcto y corrigiendo o
completando lo que no se ajustaba al material de la materia.

\newpage

%=====================================================================
%  Fragmentos (orden cronológico)
%=====================================================================
%=====================================================================
\section{Final --- 21/07/2021}
%=====================================================================

% =====================================================================
\subsection{Ejercicio 1 --- Licitación transparente con AES-GCM (Tecnociudad)}
% =====================================================================

\begin{consigna}
El consejo deliberante del partido de Tecnociudad solicitó la construcción de un sistema
distribuido que permita realizar licitaciones de forma transparente y segura. Cuando el
partido necesita contratar cualquier tipo de servicios, genera un pliego con la descripción
de lo solicitado y espera que distintos contratistas presenten ofertas. Durante esta fase,
los costos son secretos para evitar que un oferente tenga preferencia. Una vez vencido el
plazo, se publica la propuesta de cada uno, junto con el precio, y se selecciona el que
cumple con las condiciones con el menor costo.

La implementación del sistema asegura que estas propiedades ocurran utilizando criptografía.
La descripción a alto nivel es la siguiente:
\begin{itemize}
  \item Se publica la licitación con un ID único asociado.
  \item Por cada licitación, se reciben archivos cifrados con AES-GCM de parte de los
  oferentes interesados. Los archivos incluyen un documento con el ID de la licitación, la
  descripción técnica, el CUIT de la empresa y el precio.
  \item Al terminar el periodo de la licitación los oferentes comparten la clave utilizada,
  y con ello se revisan todas las ofertas, y se adjudica la ganadora.
  \item Si un oferente no provee la clave, queda automáticamente descalificado.
  \item Para transparencia, se publican todas las ofertas recibidas ya descifradas y la
  adjudicada en un portal público que puede ser consultado libremente.
\end{itemize}

\begin{enumerate}
  \item[a)] Explicar por qué ningún oferente puede alterar su propuesta sin ser detectado.
  (1 pto.)
  \item[b)] Qué medidas adicionales serían necesarias para evitar que algún administrador
  del sistema ignore a un proveedor diciendo que nunca entregó la clave. (1 pto.)
  \item[c)] Considerar el caso de un oferente que envía un archivo a nombre de un competidor,
  con un precio irrisoriamente bajo que lo perjudicaría. \textquestiondown{}Cómo evitaría esto usando funciones
  criptográficas? (1 pto.)
\end{enumerate}
\end{consigna}

\paragraph{Temas.}
Primera Parte, Clase 3 --- MACs y Cifrado Autenticado (\S Modos AEAD en la práctica: GCM;
\S Definición de cifrado autenticado; \S No-repudio: por qué los MAC no lo dan). Primera
Parte, Clase 4 --- Asimétrico y Firma Digital (\S Diferencia entre cifrar y firmar;
\S RSA-Signature / firma del hash; \S PKI y certificados digitales, no-repudio).

\paragraph{Qué necesitás saber.}
\begin{itemize}
  \item \textbf{AES-GCM} es un modo AEAD (\emph{cifrado autenticado}): además de
  confidencialidad (modo Counter), produce una etiqueta de autenticación sobre el
  \textbf{texto cifrado}. Al descifrar, primero se verifica el tag; si el cifrado fue
  modificado, la verificación \textbf{falla} y el descifrado se rechaza (no se obtiene
  ``basura'' silenciosa como con cifrado no autenticado).
  \item \textbf{Encrypt-then-MAC / AEAD} da \textbf{confidencialidad + integridad +
  autenticación}, pero \textbf{con clave simétrica compartida NO da no-repudio}: quien
  verifica el tag/GHASH conoce la misma clave que quien lo generó, así que un tercero
  no puede distinguir quién produjo el mensaje (ambas partes podrían haberlo generado).
  \item Para \textbf{no-repudio} hace falta criptografía \textbf{asimétrica}: \textbf{firmar}
  (con la clave \textbf{privada} del firmante) es lo que da no-repudio, porque solo el dueño
  de la privada pudo producir esa firma, y cualquiera la verifica con la pública sin poder
  falsificarla.
\end{itemize}

\paragraph{Notas y conexiones.}
El escenario combina dos capas que el material distingue con cuidado: (i) AES-GCM protege
la \emph{confidencialidad e integridad del archivo cifrado en tránsito/reposo} durante la
licitación (nadie puede leer los precios antes de tiempo ni modificar un archivo ya
enviado sin que se note), pero (ii) por ser \textbf{simétrico y con clave compartida}, no
resuelve el problema de que un administrador (que también conoce/gestiona las claves)
niegue haber recibido algo, ni impide que alguien se haga pasar por otro oferente. Ambos
puntos requieren la capa de \textbf{firma digital / PKI} de la Clase 4, que sí da
autenticación de origen y no-repudio.

\paragraph{Resolución.}

\textbf{a)} Ningún oferente puede alterar su propuesta sin ser detectado porque el archivo
se envía cifrado con \textbf{AES-GCM}, que es un esquema de \textbf{cifrado autenticado}
(AEAD): además de cifrar el contenido (confidencialidad), GCM calcula una \textbf{etiqueta
de autenticación} sobre el texto cifrado (usando GHASH). Al momento de descifrar (una vez
compartida la clave), se \textbf{verifica primero el tag}; si el archivo cifrado fue
modificado en cualquier bit respecto del que se envió originalmente, la verificación del
tag \textbf{falla} y el descifrado se rechaza explícitamente (a diferencia de un cifrado
sin autenticar, donde la manipulación pasaría desapercibida y solo produciría datos
``basura'' al descifrar). Por lo tanto, cualquier intento de alterar el precio o la
descripción técnica \emph{después} de enviado el archivo es detectable al momento de
abrir las ofertas.

\textbf{b)} El esquema descripto usa una \textbf{clave simétrica compartida} (el oferente
``comparte la clave utilizada''); como se remarca en la Clase 3, un esquema con clave
compartida (MAC/AEAD simétrico) \textbf{no da no-repudio}: quien controla el sistema
(el administrador) también podría --- al menos en principio --- estar en condiciones de
negar la recepción de la clave sin que el oferente pueda probar lo contrario, ya que no hay
una firma verificable por un tercero independiente. La medida adicional es introducir
\textbf{firma digital (asimétrica)}: que el oferente, al enviar la clave de descifrado,
la \textbf{firme con su clave privada}, y que el administrador, al recibirla, devuelva un
\textbf{acuse de recibo firmado} con su propia clave privada. De esta forma:
\begin{itemize}
  \item La firma del oferente sobre el envío de la clave prueba, ante cualquier tercero
  (verificando con la clave pública del oferente), que efectivamente la envió.
  \item El acuse de recibo firmado por el administrador (con su clave privada, verificable
  con su pública conocida vía PKI/certificado) es una prueba de que el sistema la recibió;
  como solo el administrador pudo generar esa firma, \textbf{no puede negarlo después}
  (no-repudio).
\end{itemize}
Esto traslada el problema de ``palabra contra palabra'' a un esquema con \textbf{propiedades
de no-repudio}, que el cifrado simétrico por sí solo no ofrece.

\textbf{c)} El problema es que el sistema, tal como está descripto, no verifica que quien
envía un archivo ``a nombre de'' un CUIT sea realmente el dueño de esa identidad: cualquiera
podría escribir el CUIT de un competidor dentro del documento cifrado. Esto se evita
usando \textbf{firma digital}: cada oferente debe \textbf{firmar su propuesta} (el
documento completo: ID de licitación, descripción técnica, CUIT y precio) con su
\textbf{clave privada} antes de cifrarlo y enviarlo, y esa clave privada debe estar
asociada a su identidad mediante un \textbf{certificado digital} emitido por una
\textbf{CA} (PKI), tal como en el ejemplo de firma de binarios de la Clase 4. Al
adjudicar, el administrador (o cualquiera con la clave pública certificada del oferente)
\textbf{verifica la firma}: si el documento dice ser del competidor pero está firmado con
la clave privada del atacante, la verificación con la clave pública \emph{del competidor}
\textbf{falla}, y la propuesta fraudulenta se descarta. En síntesis: \textbf{AES-GCM} da
confidencialidad e integridad del archivo, pero la \textbf{autenticidad del emisor} (que
el CUIT declarado sea quien realmente presentó la oferta) requiere \textbf{firma digital
con PKI}, que es justamente lo que impide la suplantación de identidad descripta.

% =====================================================================
\subsection{Ejercicio 2 --- Elección de supervisor mediante hash y XOR}
% =====================================================================

\begin{consigna}
El siguiente protocolo fue propuesto para que $N$ nodos de un sistema puedan elegir un
supervisor, de tal manera que ningún nodo pueda ganar o perder sistemáticamente
manipulando el intercambio:
\begin{enumerate}
  \item Cada nodo genera un valor $r_i$ de 64 bits y envía al resto $H(r_i)$.
  \item Luego de recibir todos los valores, cada nodo publica $r_i$.
  \item Cada nodo calcula $r_x = r_1 \oplus r_2 \oplus \dots \oplus r_n$.
\end{enumerate}
El nodo a menor distancia de $r_x$ es el nuevo supervisor, utilizando la distancia de
Hamming como métrica (se comparan bit a bit los dígitos binarios de los dos valores, y
cada discrepancia suma 1 de distancia --- por ejemplo 0010 y 1110 están a distancia 2).
En caso de empate, se reinicia.

\begin{enumerate}
  \item[a)] Explicar por qué el paso 1 es necesario y no se puede comenzar el intercambio
  en el paso 2 directamente. (1 pto.)
  \item[b)] Considerar el caso de solo tres nodos. \textquestiondown{}Qué pasaría si el nodo 2, por ejemplo,
  decide abusar el sistema, y repite en el paso 1 el valor $H(r_1)$ como si fuese el suyo?
  Proponer una modificación simple que evite el problema. (1,5 ptos.)
  \item[c)] Proponer una mejora al algoritmo para que no sea necesario repetir todo el
  intercambio en caso de empate y siempre se llegue a un resultado que no pueda ser
  manipulado. (0.5 ptos.)
\end{enumerate}
\end{consigna}

\paragraph{Temas.}
Primera Parte, Clase 3 --- MACs y Cifrado Autenticado (\S Funciones de hash: resistencia a
preimagen, segunda preimagen y colisiones; uso como \emph{compromiso}/commitment previo a
revelar un valor). \textbf{Nota:} el protocolo en sí (elección de supervisor con XOR y
distancia de Hamming) no está descripto literalmente en el material como esquema con
nombre propio; la resolución se apoya en las \textbf{propiedades de las funciones de hash}
que sí están desarrolladas en la Clase 3, aplicadas como un esquema de compromiso
(committing a $r_i$ vía $H(r_i)$ antes de revelarlo).

\paragraph{Qué necesitás saber.}
\begin{itemize}
  \item \textbf{Resistencia a preimagen:} dado un hash $y$, es computacionalmente inviable
  hallar un mensaje $x$ tal que $H(x)=y$. Este es el uso relevante acá: publicar $H(r_i)$
  primero \emph{compromete} al nodo con un valor sin revelarlo, porque nadie puede, a partir
  del hash, deducir $r_i$ para ajustar su propio valor en función de los demás.
  \item \textbf{Resistencia a segunda preimagen:} dado un mensaje $x$, es inviable hallar
  $x'\neq x$ con el mismo hash. Esto es lo que impide, una vez publicado $H(r_i)$, cambiar
  el $r_i$ real por otro $r_i'$ que también hashee a lo mismo en el paso 2.
  \item El material señala explícitamente el uso de compromisos criptográficos (``probar que
  se tenía algo antes de revelarlo''; ``juegos de cartas con compromiso'') como aplicación de
  la resistencia a (segunda) preimagen.
\end{itemize}

\paragraph{Notas y conexiones.}
El paso 1 (publicar $H(r_i)$ antes de $r_i$) es exactamente el patrón de
\textbf{compromiso-y-revelación} (commit-then-reveal) que la Clase 3 asocia a la resistencia
a preimagen/segunda preimagen: sirve para que ningún participante pueda elegir su valor
\emph{después} de ver los de los demás (lo cual rompería la imparcialidad de la elección).
Es el mismo principio que sustenta, por ejemplo, no poder generar de antemano el $r_i$ en
función del resultado final.

\paragraph{Resolución.}

\textbf{a)} El paso 1 (publicar $H(r_i)$ antes de $r_i$) es necesario porque, si en cambio
se comenzara directamente publicando los $r_i$ en claro (paso 2), un nodo malicioso podría
esperar a ver los valores $r_j$ ($j\neq i$) que ya publicaron todos los demás y \textbf{recién
entonces} elegir su propio $r_i$ de forma tal que el XOR resultante $r_x$ quede
convenientemente cerca (en distancia de Hamming) de su propio valor, garantizándose ser
supervisor. Publicar primero $H(r_i)$ actúa como un \textbf{compromiso}: por
\textbf{resistencia a preimagen}, ningún otro nodo puede, a partir de $H(r_i)$, obtener
$r_i$ antes de que se revele en el paso 2, de modo que nadie puede ajustar su propio valor
en función de los $r_i$ de los demás. Así se garantiza que todos los $r_i$ quedaron
``congelados'' antes de que ninguno conociera el de los otros, haciendo el intercambio
imparcial.

\textbf{b)} Con tres nodos, si el nodo 2 en el paso 1 publica $H(z)=H(r_1)$ (repite el
compromiso del nodo 1) en lugar de comprometerse con su propio valor, en el paso 2 puede
publicar $z=r_1$ como si fuera su $r_2$ (es un valor válido porque su hash coincide con lo
publicado). El resultado sería $r_x = r_1\oplus r_1\oplus r_3 = r_3$ (el XOR de $r_1$ consigo
mismo se cancela), y el nodo 2 conocería de antemano que $r_x=r_3$: puede entonces mirar la
distancia de Hamming entre $r_1,r_2^{\text{real}}$ y $r_3$ contra $r_x=r_3$ y decidir
manipular el resultado a su favor (por ejemplo, el propio nodo 3 quedaría siempre exactamente
a distancia 0 de $r_x$, y el nodo 2 puede jugar con esa certeza, o directamente evitar que el
nodo 1 o 3 ganen si le conviniera). En cualquier caso, se rompe la propiedad buscada de que
``ningún nodo pueda ganar o perder sistemáticamente manipulando el intercambio''.

\textbf{Modificación simple:} exigir que \textbf{todos los compromisos $H(r_i)$ publicados en
el paso 1 sean distintos entre sí} ($H(r_i)\neq H(r_j)$ para todo $i\neq j$); si dos nodos
publican el mismo hash, se detecta la colisión de compromisos y se reinicia el proceso (o se
descalifica al nodo repetido). Como por \textbf{resistencia a segunda preimagen} es inviable
para el nodo 2 encontrar un $r_2'\neq r_1$ propio con $H(r_2')=H(r_1)$, la única forma de
publicar un hash igual al de otro nodo es copiándolo literalmente, lo cual esta verificación
detecta y bloquea antes de llegar al paso 2.

\textbf{c)} Para no tener que reiniciar todo el intercambio (generar nuevos $r_i$ y repetir
los tres pasos) en caso de empate en la distancia de Hamming, se puede definir un
\textbf{criterio de desempate determinístico y no manipulable} sobre la información que ya
está publicada, sin generar nueva aleatoriedad: por ejemplo, entre los nodos empatados en
distancia mínima a $r_x$, elegir como supervisor al que tenga \textbf{el valor $r_i$
numéricamente mayor (o menor)} interpretado como entero, o el que primero coincide bit a bit
con $r_x$ leyendo de izquierda a derecha. Como los $r_i$ ya fueron comprometidos con hash y
revelados (y no pueden cambiarse sin romper la resistencia a segunda preimagen del
compromiso), este criterio de desempate no puede ser manipulado por ningún nodo a esta
altura del protocolo (todos los valores ya están fijados e inmutables), y siempre produce un
único ganador sin necesidad de reiniciar el intercambio completo.

% =====================================================================
\subsection{Ejercicio 3 --- Pentesting con hipótesis de falla (Agrolin)}
% =====================================================================

\begin{consigna}
Agrolin S.A.\ se dedica a la construcción de soluciones tecnológicas para el campo.

Recientemente se publicó un informe que menciona múltiples vulnerabilidades en uno de sus
productos y provocó muchas pérdidas de ventas a punto de concretarse. Para corregir la
situación, la empresa decidió contratar un servicio de \emph{pen testing}, realizar una
evaluación más exhaustiva y corregir el producto para empezar a recuperar el daño a su
imagen.

El producto en cuestión permite hacer un seguimiento del rendimiento de las cosechas a lo
largo del tiempo y determinar las rotaciones óptimas para mejorar la productividad. El
sistema está compuesto por un frontend web que es un \emph{single page web app}, donde el
usuario puede registrarse, autenticarse, dar de alta lotes marcado sus límites en un mapa,
actualizarlos o eliminarlos. El sistema muestra según datos históricos de todos los clientes
de la zona, el rinde esperado para la superficie en cuestión, y sugiere el tipo de cultivo
con mayor rentabilidad en función de los valores de rinde y precios de mercados futuros
estimados, y permite que el usuario ingrese los datos del tipo de cultivo y luego de la
cosecha el rinde obtenido para comparar y poder mejorar en las próximas sugerencias.

La aplicación web se ejecuta en los navegadores de cada cliente, y hace llamadas a un
servidor backend a través de una API REST.

Siguiendo la metodología de hipótesis de falla, establecer 4 hipótesis de vulnerabilidades
que tengan alta probabilidad de existir en el escenario descripto. Para cada una, describir
únicamente la hipótesis y la prueba que realizaría para confirmarla o refutarla (1 pto.\ por
hipótesis).

\emph{Aclaraciones:} No se considerarán válidas hipótesis genéricas. Tampoco se considerarán
dos hipótesis que sean variantes del mismo tipo de falla (por ejemplo, el mismo problema en
dos pantallas distintas).
\end{consigna}

\paragraph{Temas.}
Segunda Parte, Clase 6 --- Pentesting (\S Metodología de Hipótesis de Falla: los pasos;
\S Paso 2 --- Planteo de hipótesis de falla). Segunda Parte, Clase 3 --- Principios de
Diseño y Vulnerabilidades (\S CWE / Top 25 MITRE: SQL Injection CWE-89; \S OWASP Top 10;
\S STRIDE: Denial of Service, Spoofing/Improper Authentication).

\paragraph{Qué necesitás saber.}
\begin{itemize}
  \item La \textbf{metodología de hipótesis de falla} exige, en el paso de
  \textbf{planteo de hipótesis}, asumir la existencia de una vulnerabilidad concreta
  \emph{antes} de probarla; luego se define una \textbf{prueba} específica para
  confirmarla o refutarla (no alcanza con nombrar la categoría genérica: hay que decir
  \emph{qué} se sospecha y \emph{cómo} se lo comprueba en este sistema puntual).
  \item El material lista, dentro del \textbf{Top 25 CWE de MITRE}, entre las debilidades
  más comunes: \textbf{CWE-89 SQL Injection} (falta de control en el ingreso de datos que
  permite ejecutar comandos SQL: los datos del usuario se mezclan sin sanitizar con el
  comando y la base los interpreta como código; defensa: \emph{prepared statements}/
  parametrizar), \textbf{CWE-79 Cross-site Scripting}, y (Clase 3, taxonomía STRIDE)
  \textbf{Denial of Service} (agotar los recursos necesarios para dar el servicio,
  afectando disponibilidad) como categoría de amenaza reconocida explícitamente como
  \emph{ataque de seguridad} aunque no comprometa confidencialidad/integridad.
  \item \textbf{OWASP} publica el Top 10 de vulnerabilidades web más comunes (frameworks de
  referencia junto con MITRE), relevante para un frontend \emph{single page web app} +
  API REST como el descripto (autenticación, autorización, control de acceso a recursos por
  ID).
\end{itemize}

\paragraph{Notas y conexiones.}
El enunciado da todas las pistas típicas de una aplicación web con API REST y autenticación
de usuarios: (1) el sistema tiene \textbf{autenticación} (registrarse/autenticarse), (2)
tiene \textbf{ingreso de datos de usuario} que probablemente lleguen a una base de datos
(lotes, límites en mapa, tipo de cultivo, rinde), (3) hace \textbf{consultas agregadas
costosas} sobre datos históricos de todos los clientes de la zona, y (4) maneja
\textbf{recursos por identificador} (lotes) que pertenecen a distintos clientes. Cada una
de estas pistas mapea a una categoría de vulnerabilidad del Top 25/CWE o de STRIDE vista en
la Clase 3, y la consigna pide aplicar sobre ellas el paso de \textbf{hipótesis + prueba}
de la Clase 6 (no basta con nombrar la categoría).

\paragraph{Resolución.}

\textbf{Hipótesis 1 --- Autenticación deficiente (Improper Authentication).}
\emph{Hipótesis:} el sistema de autenticación del frontend/API es débil o directamente
inexistente en algunos endpoints, permitiendo iniciar sesión con credenciales triviales o
acceder a funciones que deberían requerir sesión válida sin ella.
\emph{Prueba:} intentar registrarse con credenciales simples (usuario/contraseña
predecibles) y verificar si el sistema las acepta sin exigir complejidad; adicionalmente,
intentar invocar directamente los endpoints de la API REST (alta/consulta de lotes) sin
enviar token/cookie de sesión, para verificar si el backend responde igual que a un usuario
autenticado.

\textbf{Hipótesis 2 --- SQL Injection (CWE-89) en el ingreso de datos.}
\emph{Hipótesis:} los campos donde el usuario ingresa datos (tipo de cultivo, límites del
lote, rinde obtenido) no sanitizan/parametrizan la entrada antes de construir la consulta
a la base de datos, permitiendo inyectar comandos SQL.
\emph{Prueba:} en un campo de texto (por ejemplo, nombre del cultivo), ingresar una cadena
con sintaxis SQL (p.\ ej.\ un valor que cierre la comilla e intente insertar/alterar una
tabla) y verificar, sin necesidad de dañar la base, si el sistema devuelve un error de
sintaxis SQL expuesto o un comportamiento anómalo que indique que el dato fue interpretado
como comando en lugar de como valor.

\textbf{Hipótesis 3 --- Denial of Service por consultas agregadas costosas.}
\emph{Hipótesis:} la funcionalidad que muestra, según datos históricos de \textbf{todos los
clientes de la zona}, el rinde esperado y sugiere el cultivo más rentable en función de
precios de mercados futuros, requiere procesar/parsear múltiples consultas pesadas por cada
solicitud; enviar varias de estas solicitudes en paralelo podría agotar los recursos del
servidor y denegar el servicio a otros usuarios.
\emph{Prueba (acotada, sin ejecutar el ataque real):} verificar únicamente si es posible
disparar 2 de estas consultas costosas en paralelo desde una misma sesión/IP sin que el
sistema las limite (rate limiting) o las rechace; no se realiza el ataque a mayor escala
para no afectar el servicio, y se reporta la situación si se confirma la ausencia de
límites.

\textbf{Hipótesis 4 --- Autorización deficiente / control de acceso a recursos de otro
cliente (Improper Authorization / Broken Access Control).}
\emph{Hipótesis:} una vez autenticado, el sistema no verifica correctamente que el usuario
solo pueda acceder a sus propios lotes, permitiendo acceder a lotes (y sus límites en el
mapa, rindes, etc.) de otros clientes modificando un identificador en la URL o en la
llamada a la API REST.
\emph{Prueba:} tras crear una cuenta y autenticarse, identificar la URL/endpoint que expone
un lote propio con su identificador (p.\ ej.\ \texttt{.../cliente/12/campo/1}) y probar
acceder al mismo recurso cambiando el identificador de cliente por uno ajeno (p.\ ej.\
\texttt{.../cliente/1/campo/1}); si el sistema devuelve los datos sin validar que el lote
pertenece al usuario autenticado, se confirma la hipótesis. No se modifican datos ajenos
para no ser intrusivo; solo se confirma el acceso indebido.

\textbf{Nota sobre alcance:} las cuatro hipótesis corresponden a categorías distintas de
falla (autenticación, inyección de datos, disponibilidad y autorización), tal como exige la
aclaración de la consigna, y todas están sustentadas en la taxonomía CWE/STRIDE/OWASP
desarrollada en la Clase 3 y en el paso de ``planteo de hipótesis'' de la metodología de
pentesting de la Clase 6.

\clearpage
%=====================================================================
\section{Final --- 17/07/2023}
%=====================================================================

\subsection{Ejercicio 1 --- Dispositivos de medición de agua (AES-ECB, HTTPS, integridad)}

\begin{consigna}
Hay una empresa que se encarga de controlar el agua. Para esto la empresa implementa unos
dispositivos encargados de hacer las respectivas mediciones. Los mismos le informan a la
central los datos recolectados.

El procedimiento cuando se pone un dispositivo es el siguiente: primero el dispositivo emite
su ID por broadcast, luego la central responde a ese ID y le envía una llave por HTTP. Cuando
el dispositivo le debe enviar lo recolectado a la empresa usa AES-ECB como método de
encripción.

La empresa contrata una auditoría y le dice que posee fallas en la comunicación de la llave y
problemas de integridad en la comunicación de los dispositivos.
\begin{enumerate}
  \item[a)] Usando HTTPS \textquestiondown{}se hubieran solucionado todos los problemas que informó la
  auditoría?
  \item[b)] En caso de no poder usar HTTPS por cuestiones de presupuestos \textquestiondown{}qué otro protocolo
  utilizaría para el intercambio de claves?
  \item[c)] \textquestiondown{}Qué cambios son necesarios para asegurar la integridad en el envío de información
  de los dispositivos hacia la central?
\end{enumerate}
\end{consigna}

\paragraph{Temas.}
Primera Parte, Clase 2 --- Cifrado (\S ECB, \S modos de cifrado, \S CPA-security); Clase 3 ---
MACs y Cifrado Autenticado (\S qué es un MAC, \S cifrado autenticado, \S modos AEAD:
CCM/GCM); Clase 5 --- Protocolos (\S TLS Handshake, \S Diffie--Hellman); Segunda Parte,
Clase 11 --- Protección de Datos (\S protocolos de transporte cifrados: TLS/HTTPS).

\paragraph{Qué necesitás saber.}
\begin{itemize}
  \item \textbf{ECB (Electronic Codebook):} cifra cada bloque de forma independiente,
  $c_i=e_k(m_i)$. Es determinista: bloques de texto plano iguales producen cifrados iguales
  $\Rightarrow$ \textbf{revela patrones/estructura} y \textbf{no es CPA-Secure}. No debe
  usarse.
  \item \textbf{HTTPS = HTTP sobre TLS.} El TLS Handshake acuerda un método de intercambio de
  clave (RSA, Diffie--Hellman con certificado, ECDH, etc.), autentica con firmas y deriva un
  \emph{Master Secret}; la criptografía negociable para el resto de la sesión incluye cifrado
  simétrico (recomendado AES-GCM) y HMAC/AEAD para integridad.
  \item \textbf{Diffie--Hellman} (Clase 5): acuerdo de clave cuya seguridad descansa en el
  problema del logaritmo discreto (DLP); permite que ambas partes compartan una clave secreta
  sobre un canal inseguro sin haberla intercambiado antes. \textbf{DH puro no autentica}: es
  vulnerable a Man-in-the-Middle si no se combina con certificados (por eso en TLS va con
  RSA/DSS).
  \item \textbf{Cifrado autenticado / AEAD} (Clase 3): en la práctica, si se necesita
  confidencialidad casi siempre también se necesita integridad, porque el cifrado simple
  (como ECB) puede ser manipulado sin que nadie lo note. La recomendación es usar un modo
  \textbf{AEAD} como \textbf{AES-GCM} (o CCM), que da confidencialidad e integridad juntas.
\end{itemize}

\paragraph{Notas y conexiones.}
El escenario combina dos fallas de diseño típicas de la materia: (1) enviar una clave en
claro por un canal no cifrado (HTTP) para luego usarla en un cifrado simétrico, y (2) usar
ECB, el modo que la Clase 2 marca como ``NO utilizar'' por no ser CPA-Secure y no dar
integridad. Conecta con la distinción Clase 3 entre cifrado (confidencialidad) y MAC/AEAD
(integridad): un cifrado por sí solo, aunque fuera un modo seguro como CBC o GCM sin
autenticar, no impide que un atacante manipule bits del cifrado sin ser detectado si no hay
un mecanismo de integridad de por medio.

\paragraph{Resolución.}

\textbf{a) No, no necesariamente se hubieran solucionado todos los problemas.}
\begin{itemize}
  \item Usando \textbf{HTTPS con TLS} para el paso donde la central envía la llave, sí se
  soluciona el problema de la \textbf{comunicación de la llave}: TLS resuelve el intercambio
  de claves de forma autenticada (certificados) y cifra el canal, evitando que un atacante
  intercepte la llave en texto plano durante el handshake inicial. Esto ataca directamente la
  falla ``fallas en la comunicación de la llave''.
  \item Sin embargo, HTTPS solo protege la \textbf{comunicación} del dispositivo con la
  central (la llave viajando por HTTP). \textbf{No cambia el método de encripción que el
  dispositivo usa para enviar lo recolectado}, que sigue siendo \textbf{AES-ECB}. Como ECB es
  determinista y no es CPA-Secure, \textbf{no da integridad} (bloques de datos repetidos
  producen cifrados idénticos, y un atacante puede reordenar/reemplazar bloques de cifrado
  sin que el receptor lo detecte, ya que ECB no incluye ningún mecanismo de autenticación).
  Por lo tanto el problema de \textbf{integridad en la comunicación de los dispositivos}
  \textbf{no queda resuelto} solo con HTTPS en el paso de intercambio de la llave.
\end{itemize}

\textbf{b) Diffie--Hellman.} Si no se puede pagar el costo de HTTPS/TLS completo, se puede
usar directamente un protocolo de \textbf{intercambio de clave Diffie--Hellman}: permite que
la central y el dispositivo acuerden una clave secreta compartida sobre un canal inseguro sin
haberla compartido antes, basando su seguridad en el problema del logaritmo discreto. Queda
como salvedad (marcada en el material) que \textbf{DH puro no autentica a las partes} y es
susceptible a un ataque Man-in-the-Middle si un atacante se interpone en el intercambio; para
mitigarlo haría falta combinarlo con algún mecanismo de autenticación (certificados o claves
públicas ya conocidas), tal como TLS lo hace con RSA/DSS.

\textbf{c)} Para asegurar la integridad en el envío de la información recolectada hacia la
central hace falta:
\begin{enumerate}
  \item \textbf{Dejar de usar ECB en solitario:} es un modo sin ningún mecanismo de
  autenticación, determinista y no CPA-Secure.
  \item \textbf{Pasar a un esquema de cifrado autenticado (AEAD)}, que da confidencialidad
  \emph{e} integridad en una sola construcción: \textbf{AES-GCM} (recomendación práctica del
  material: rápido, estandarizado, en todas las librerías serias) o \textbf{AES-CCM}. Ambos
  agregan, además del cifrado, un \emph{tag} de autenticación que permite al receptor
  verificar que el mensaje cifrado no fue modificado en el camino.
  \item Alternativamente (o en conjunto), agregar explícitamente un \textbf{MAC/HMAC} sobre lo
  enviado, de forma que la central pueda \textbf{chequear que el mensaje no fue modificado en
  el camino} (verificar el tag con la clave compartida antes de aceptar la medición).
\end{enumerate}
Nota de fidelidad: el material no detalla un protocolo específico de bajo costo distinto de
Diffie--Hellman para IoT (p.\ ej. protocolos livianos de sensores); la respuesta se limita a
lo cubierto en las clases de Protocolos y MACs/Cifrado Autenticado.

%=====================================================================
\subsection{Ejercicio 2 --- Ocultar un mensaje ``tipo sombras'' (OTP + integridad con hash/HMAC)}
%=====================================================================

\begin{consigna}
Se quiere ocultar un mensaje de forma segura con una técnica similar a las sombras. Dado un
mensaje ``m'' se eligen dos números random ($x_1,x_2$) y se generan por dos canales separados
(posiblemente en lugares distintos):
\begin{itemize}
  \item $p_1 = x_1 \oplus m \oplus x_2$
  \item $p_2 = x_2 \oplus x_1$
  \item $p_1$ y $h(p_2)$
  \item $p_2$ y $h(p_1)$
\end{itemize}
\begin{enumerate}
  \item[a)] Explicar por qué si un atacante tiene acceso a solo 1 canal de comunicación aun se
  mantiene la confidencialidad.
  \item[b)] Explicar por qué si un atacante tiene acceso a un solo canal de comunicación aun se
  mantiene la integridad.
  \item[c)] Si un atacante tiene acceso a los 2 canales de información, cómo hacer para que se
  mantenga \underline{solo la integridad}.
\end{enumerate}
\end{consigna}

\paragraph{Temas.}
Primera Parte, Clase 2 --- Cifrado (\S secreto perfecto, \S One-Time Pad / Vernam); Clase 5 ---
Protocolos (\S secreto compartido de Shamir, umbral $(t,n)$: con menos del umbral no se filtra
información); Clase 3 --- MACs y Cifrado Autenticado (\S funciones de hash, \S HMAC,
\S no-repudio: por qué MAC/hash no lo dan).

\paragraph{Qué necesitás saber.}
\begin{itemize}
  \item \textbf{OTP (Clase 2):} $e_k(m) = m \oplus k$ con $k$ uniforme y del mismo tamaño que
  $m$ tiene \textbf{secreto perfecto}: $\Pr[M{=}m\mid C{=}c] = \Pr[M{=}m]$. Si $x_1$ y $x_2$
  son aleatorios e independientes, $x_1\oplus x_2$ se comporta como una clave uniforme.
  \item \textbf{Secreto compartido / esquema umbral (Clase 5):} repartir información en
  ``sombras'' tal que \textbf{con menos de las necesarias no se filtra ninguna información}
  del secreto (entropía máxima). Acá el mensaje se reparte en dos partes ($p_1,p_2$) tal que
  ninguna por separado revela $m$.
  \item \textbf{Función de hash (Clase 3):} \textbf{no tiene clave} (es información pública);
  depende de todos los bits del mensaje (efecto avalancha) y permite detectar si el mensaje
  fue alterado comparando el hash recibido contra el hash recalculado. Al no tener clave, un
  atacante que \textbf{controla el mensaje y puede recalcular el hash él mismo} puede
  modificar ambos de forma consistente sin ser detectado.
  \item \textbf{HMAC (Clase 3):} MAC construido a partir de una función de hash y una
  \textbf{clave} $k$ compartida; a diferencia del hash simple, el atacante sin la clave
  \textbf{no puede recalcular un tag válido} aunque conozca y modifique el mensaje.
\end{itemize}

\paragraph{Notas y conexiones.}
El esquema es análogo a un \textbf{secreto compartido de 2 sombras} (Clase 5) combinado con
\textbf{OTP} (Clase 2): $p_1$ es el mensaje enmascarado con el XOR de dos valores aleatorios
(equivalente a una clave OTP $x_1\oplus x_2$), y $p_2$ es esa misma máscara sin el mensaje.
Ninguna de las dos partes por sí sola alcanza para reconstruir $m$: hace falta juntar ambos
canales, igual que con el umbral de Shamir se necesita juntar $t$ sombras. El agregado de
$h(\cdot)$ cruzado conecta con Clase 3 (funciones de hash sin clave) y anticipa la necesidad
de HMAC para integridad \emph{autenticada} cuando el atacante controla todo.

\paragraph{Resolución.}

\textbf{a) Confidencialidad con acceso a un solo canal.}
\begin{itemize}
  \item \textbf{Canal 1} entrega $(p_1, h(p_2))$, con $p_1 = x_1\oplus m\oplus x_2$. Como
  $x_1$ y $x_2$ son elegidos al azar e independientes, $k = x_1\oplus x_2$ se comporta como
  una \textbf{clave OTP uniforme}: $p_1 = m \oplus k$. Por el resultado de secreto perfecto
  del OTP, sin conocer $k$ el atacante no obtiene ninguna información sobre $m$ a partir de
  $p_1$ solo. $h(p_2)$ tampoco aporta nada sobre $m$ (depende solo de $p_2 = x_1\oplus x_2$,
  que no involucra a $m$).
  \item \textbf{Canal 2} entrega $(p_2, h(p_1))$, con $p_2 = x_1\oplus x_2$. Esto es
  exactamente la máscara $k$ usada para enmascarar $m$, pero \textbf{sin el mensaje}: no
  contiene ninguna información de $m$. $h(p_1)$ tampoco revela $m$ porque una función de hash
  no es invertible (no permite recuperar el mensaje original a partir del resumen).
\end{itemize}
En ambos casos se necesitan \textbf{ambos canales} para calcular
$p_1 \oplus p_2 = (x_1\oplus m\oplus x_2)\oplus(x_1\oplus x_2) = m$; con uno solo la
confidencialidad de $m$ se mantiene.

\textbf{b) Integridad con acceso a un solo canal.}
Cada canal lleva el valor a transmitir junto con el \textbf{hash del otro} valor
($p_1$ con $h(p_2)$, y $p_2$ con $h(p_1)$), es decir, hashes \textbf{cruzados}. Si el
atacante solo controla un canal (por ejemplo el que lleva $p_1$ y $h(p_2)$) y modifica $p_1$
en tránsito, no tiene forma de recalcular un $h(p_2)$ consistente porque \textbf{no tiene
acceso a $p_2$} (viaja por el otro canal, que no controla). Al llegar los dos valores al
receptor, este recalcula $h(p_2)$ recibido en el otro canal y $h(p_1)$ recibido en este, y
si $p_1$ fue alterado, el hash recalculado de $p_1$ no coincidirá con el $h(p_1)$ que viajó
por el canal 2, delatando la alteración. Con un solo canal comprometido, la verificación
cruzada contra el otro canal (no comprometido) detecta la manipulación.

\textbf{c) Con acceso a los dos canales, mantener solo la integridad.}
Si el atacante controla ambos canales, puede leer $p_1$ y $p_2$ enteros, calcular
$h(\cdot)$ él mismo (la función de hash \textbf{no usa clave}, es pública) y modificar
$p_1$ y/o $p_2$ recalculando los hashes cruzados de forma consistente: el esquema de hash
simple deja de servir para integridad en ese caso (y de hecho la confidencialidad de $m$
también se pierde, porque con los dos canales se recupera $m = p_1\oplus p_2$; por eso la
consigna pide mantener \emph{solo} la integridad, asumiendo que la confidencialidad ya no es
sostenible con acceso total a ambos canales).

Para mantener la integridad hay que reemplazar el hash simple $h(\cdot)$ por un mecanismo
\textbf{que dependa de una clave que el atacante no tenga}:
\begin{itemize}
  \item Usar \textbf{HMAC} en lugar de un hash sin clave: $\mathrm{HMAC}_k(p_2)$ en el canal 1
  y $\mathrm{HMAC}_k(p_1)$ en el canal 2, con una clave $k$ secreta compartida entre emisor y
  receptor mediante un canal previo seguro. Sin conocer $k$, el atacante no puede generar un
  tag $\mathrm{HMAC}_k$ válido para un $p_1$ o $p_2$ modificado, aunque vea y edite el
  contenido de ambos canales.
  \item Alternativamente, \textbf{firmar digitalmente} $p_1$ y $p_2$ con la clave
  \textbf{privada} del emisor: el atacante, aun controlando ambos canales, no puede falsificar
  la firma sin esa clave privada, por lo que cualquier alteración se detecta al verificar la
  firma con la clave pública correspondiente. (Esta variante, a diferencia de HMAC, además
  daría no-repudio, ya que el material remarca que un MAC/HMAC simétrico \textbf{no} da
  no-repudio porque ambas partes comparten la misma clave.)
\end{itemize}

%=====================================================================
\subsection{Ejercicio 3 --- Pentesting a una página de noticias (hipótesis de falla)}
%=====================================================================

\begin{consigna}
Una página de noticias en la que hay un usuario y un generador de noticias. A cada noticia se
le puede poner etiquetas como categorías (Deportes, Interés general) o Ubicación (Argentina).

Cada usuario lleva un puntaje (+1 por cada noticia rateada positivamente y -1 por cada noticia
rateada negativamente). La página tiene una API REST y los llamados a la misma se hacen por
HTTPS.

Se le pagará por visualizaciones al generador de la noticia teniendo en cuenta algunos
factores que no me acuerdo bien y no venían al caso.

Se hizo un análisis de la página y se encontró que:
\begin{enumerate}
  \item Existe una URL para pedir todas las noticias:
  \texttt{https://dominio.com/news?limit=100\&page=20}
  \item Cuando se le recomienda o no recomienda una noticia se hace a través de
  \texttt{https://dominio.com/\{generator\_id\}/\{news\_id\}/upvoted} o
  \texttt{https://dominio.com/\{generator\_id\}/\{news\_id\}/downvoted}
  \item Las personas se registran por usuario y password
  \item Se chequearon los inputs de la página y cualquier intento de tag HTML ingresado
  aparecerá como texto en la página $\rightarrow$ XSS no probable.
\end{enumerate}
Establecer 4 hipótesis de falla con las respectivas pruebas que haría.
\end{consigna}

\paragraph{Temas.}
Segunda Parte, Clase 10b --- Pentesting (\S Metodología de Hipótesis de Falla: los pasos en
orden --- Recolección, Hipótesis, Prueba, Generalización, Eliminación; \S Injection flaws;
\S OWASP); Clase 8 --- Principios de Diseño y Vulnerabilidades (relacionado: validación de
entrada, XSS/CWE-79).

\paragraph{Qué necesitás saber.}
\begin{itemize}
  \item La \textbf{metodología de hipótesis de falla} (Clase 10b) tiene 5 pasos en orden:
  \textbf{Recolección de información} $\to$ \textbf{Planteo de hipótesis de falla} $\to$
  \textbf{Prueba} $\to$ \textbf{Generalización} $\to$ \textbf{Eliminación} (opcional). Cada
  hipótesis debe: (1) partir de la información recolectada del sistema (los 4 hallazgos del
  enunciado), (2) suponer una vulnerabilidad concreta, y (3) tener una \textbf{prueba}
  asociada que la verifique.
  \item El material define el \textbf{XSS (CWE-79)} como inyección de JS ejecutado por el
  navegador de la víctima, y aquí el enunciado ya lo descarta explícitamente (los tags HTML
  se muestran como texto, no se ejecutan) --- es un hallazgo del paso de Recolección que cierra
  esa hipótesis.
  \item El material menciona en general \textbf{OWASP} como organización que publica el
  Top 10 de vulnerabilidades web y menciona \textbf{Injection flaws} (SQLi, XSS, OS command)
  como la familia más frecuente ($\sim$90\% según el profesor).
\end{itemize}

\begin{warnbox}{Aviso de fidelidad al material}
El enunciado del alumno (usado solo como referencia) etiqueta las hipótesis con nombres de
categorías estilo OWASP Top 10 API (``Improper Authentication'', ``Broken Access Control'',
``BOLA'', ataques de \emph{replay}). Los resúmenes de la materia \textbf{sí} cubren OWASP
como organización/Top 10 y la familia de \emph{Injection flaws}, y \textbf{sí} cubren
conceptos de control de acceso/autenticación de otras clases, pero \textbf{no} desarrollan
esas siglas puntuales (BOLA, Improper Authentication, Broken Access Control como categorías
nombradas) como parte del temario de Pentesting. Por eso, en la resolución de abajo se
plantean las 4 hipótesis con el vocabulario y la metodología efectivamente vistos en la
Clase 10b (hipótesis de falla + prueba), evitando nombrar categorías OWASP específicas que no
están en el material, y describiendo la falla en términos genéricos de control de acceso /
autenticación / lógica de negocio ya cubiertos en Segunda Parte.
\end{warnbox}

\paragraph{Resolución.}
Aplicando la metodología de hipótesis de falla (Clase 10b) sobre los 4 hallazgos de la
Recolección de información:

\begin{enumerate}
  \item \textbf{Hipótesis 1 --- falta de límite de intentos en el login.} El sistema se
  registra y autentica por usuario y contraseña (hallazgo 3), y no hay ninguna mención de
  control sobre los intentos de autenticación. \textbf{Hipótesis:} el endpoint de login no
  limita ni registra intentos fallidos, permitiendo probar contraseñas por fuerza bruta contra
  una cuenta válida. \textbf{Prueba:} realizar múltiples intentos de autenticación contra un
  usuario existente con contraseñas distintas y observar si el sistema bloquea, demora o
  registra los intentos repetidos, o si los permite indefinidamente.

  \item \textbf{Hipótesis 2 --- falta de verificación del vínculo usuario/\texttt{generator\_id}
  al votar.} Los endpoints de recomendación usan \texttt{\{generator\_id\}/\{news\_id\}/upvoted}
  y \texttt{downvoted} directamente en la URL (hallazgo 2), sin mostrar cómo se valida contra
  quién está autenticado. \textbf{Hipótesis:} el back-end no verifica que el
  \texttt{generator\_id} de la URL corresponda al usuario autenticado o autorizado para votar,
  permitiendo que un usuario vote (sume/reste puntaje) en nombre de otro generador o sin estar
  autenticado. \textbf{Prueba:} (i) sin loguearse, llamar directamente a la URL de
  \texttt{upvoted}/\texttt{downvoted} con un \texttt{generator\_id}/\texttt{news\_id} válidos y
  observar si el sistema restringe el acceso; (ii) logueado como un usuario, llamar a la misma
  URL con un \texttt{generator\_id} que no es el propio y observar si el sistema lo permite.

  \item \textbf{Hipótesis 3 --- manipulación del parámetro \texttt{limit} de la URL de
  noticias.} El hallazgo 1 expone una URL con parámetros \texttt{limit} y \texttt{page}
  controlados por el cliente. \textbf{Hipótesis:} el valor de \texttt{limit} no está validado
  ni acotado en el servidor, permitiendo pedir una cantidad arbitrariamente grande de noticias
  en una sola request (degradación del servicio / consumo excesivo de recursos).
  \textbf{Prueba:} modificar el valor de \texttt{limit} a números extremadamente altos (o
  negativos/no numéricos) y observar la respuesta y el tiempo/recursos que consume el
  servidor.

  \item \textbf{Hipótesis 4 --- falta de control sobre votos repetidos del mismo usuario.} El
  puntaje del generador depende de +1/-1 por cada rating (hallazgo del enunciado principal).
  \textbf{Hipótesis:} el sistema no valida que un mismo usuario ya haya votado una misma
  noticia, permitiendo repetir el llamado a \texttt{upvoted}/\texttt{downvoted} para inflar o
  deflar artificialmente el puntaje de un generador (equivalente a un ataque de repetición
  sobre la misma acción). \textbf{Prueba:} repetir la misma llamada de
  \texttt{upvoted}/\texttt{downvoted} (mismo usuario, misma noticia) dos o más veces y verificar
  si el puntaje del generador se incrementa cada vez o si el sistema detecta y descarta el
  duplicado.
\end{enumerate}

Nota: el hallazgo 4 del enunciado (chequeo de tags HTML) queda \textbf{descartado} como
hipótesis de XSS por el propio enunciado (paso de Recolección ya cerrado): cualquier tag
ingresado se muestra como texto, no se ejecuta, así que no aplica la familia de
\emph{Injection flaws} tipo XSS sobre ese input en particular.

\clearpage
%=====================================================================
\section{Final --- 06/07/2023}
%=====================================================================

% =====================================================================
\subsection{Ejercicio 1 --- $k$-anonimato en base de datos de emails (SPAM)}
% =====================================================================

\begin{consigna}
Considerar una base de datos de emails que han sido utilizados para publicidad no
deseada (SPAM) y un servicio que permite a cualquier persona consultar si su email está
en la lista. La consulta misma le da al dueño del sistema acceso a nuevos emails que
podrían ser vendidos, robados, o utilizados sin consentimiento.

El servicio pretende mantener el anonimato de la información, en el sentido de no poder
asociar un email particular a un usuario que usa el mismo. Para ello se decidió
implementar un mecanismo conocido como $k$-annonymity, que establece la imposibilidad de
asociar un dato a una entidad discernible en un grupo de al menos $k$ entidades.

El servicio guarda en su base de datos el resultado de aplicar un hash criptográfico
SHA-256 sobre cada email. Un cliente que quiera hacer una búsqueda, calcula
\underline{localmente} el hash del email a consultar y envía \textbf{solo los primeros
$x$ bits} del hash (donde $x$ es una constante definida como parte del servicio),
recibiendo como respuesta todos los hashes de la base de datos que tengan dicho
prefijo. De esta manera el usuario puede buscar si aparece su hash completo en la
respuesta y responder la pregunta.

\begin{enumerate}
  \item[a)] ¿Mejora la seguridad si se utiliza un SALT de 256 bits para almacenar los
  hashes? ¿Tiene alguna otra consecuencia? Aclaración: seguridad en este contexto
  refiere al modelo a anonimato deseado. (1 pto)
  \item[b)] Explicar por qué alguien que controle el servidor no puede saber ni cuál era
  el email sobre el que preguntó el usuario, ni si el resultado de la búsqueda fue
  exitoso o no. (1 pto)
  \item[c)] ¿Con qué criterio se puede definir el valor $x$ (la cantidad de bits de
  prefijo que debe enviar el cliente) si queremos garantizar 256-annonimity? (1 pto)
\end{enumerate}
\end{consigna}

\paragraph{Temas.}
Primera Parte, Clase 3 --- MACs y funciones de hash (\S Funciones de hash; \S
Resistencia a preimagen/segunda preimagen/colisión; \S Paradoja del cumpleaños).
Segunda Parte, Clase 2 --- Autenticación (\S Salt: qué hace y qué no hace un salt).
\textbf{Aclaración de fidelidad:} el término ``$k$-anonimato'' / $k$-anonymity
\textbf{no aparece} como tal en los resúmenes de la materia (Primera ni Segunda Parte);
la definición usada acá es la que trae el propio enunciado del examen. La resolución se
arma con las propiedades de funciones de hash y de salt que sí están en el material.

\paragraph{Qué necesitás saber.}
\begin{itemize}
  \item \textbf{Resistencia a preimagen:} dado un hash $y$, es computacionalmente
  imposible hallar $x$ con $h(x)=y$. Por eso el servidor, recibiendo un prefijo del
  hash, no puede ``invertirlo'' para recuperar el email en claro.
  \item \textbf{Salt:} perturbación aleatoria embebida en el hash, definida (según la
  materia) para el escenario de \emph{cracking offline de contraseñas}: no agranda el
  espacio de claves del usuario, pero multiplica por $2^{\text{bits de salt}}$ el costo
  del atacante que intenta adivinar el valor original, porque impide reutilizar un
  cómputo $h(a)$ contra varias cuentas en paralelo (cada cuenta usa una perturbación
  distinta).
  \item \textbf{Efecto avalancha / tamaño de salida fijo:} el hash de un mensaje
  arbitrario da una etiqueta de tamaño fijo (acá 256 bits); dos entradas distintas dan
  salidas que no guardan relación aparente entre sí.
\end{itemize}

\paragraph{Notas y conexiones.}
El protocolo del enunciado es, en esencia, un esquema de \textbf{PIR (búsqueda privada)
construido ad-hoc} sobre resistencia a preimagen: el cliente nunca envía el email en
claro, solo un prefijo de su hash, y el servidor responde con \emph{todos} los hashes
que matchean ese prefijo (no solo el buscado) para no revelar cuál interesa
puntualmente. Esto conecta con la idea de la materia de que un hash \textbf{no es
reversible} (``un hash sirve para encriptar'' es FALSO, unidad de MACs) y con la
distinción salt vs.\ espacio de claves (unidad de Autenticación): acá el salt no
protege contra fuerza bruta sobre el hash sino que \textbf{rompe el protocolo mismo},
como se ve en el inciso a).

\paragraph{Resolución.}

\textbf{a)} \textbf{No mejora la seguridad/anonimato buscado; al contrario, rompe el
protocolo.} El objetivo del esquema no es dificultar un ataque de fuerza bruta offline
contra el hash (que es para lo que el salt sirve en el material: penalizar al atacante
que intenta adivinar $x$ tal que $h(x)=y$), sino permitir que el propio usuario legítimo
recompute \emph{localmente} el hash de su email para poder consultarlo. Si se agrega un
salt $s$ por registro ($h(\text{email}\|s)$), el cliente que quiere consultar
\textbf{no conoce el salt} usado para su propio email en la base (el salt es
justamente aleatorio y --- para que sirva su propósito de la materia --- no se le
entrega al usuario de antemano), por lo que no puede recalcular localmente el prefijo
correcto para preguntar. Consecuencias:
\begin{itemize}
  \item Si el servicio de todos modos le entregara el salt correspondiente a su email
  para que pueda calcular el hash, eso \textbf{ya revela al servidor qué salt se pidió},
  y si el salt es por-email, pedir ``el salt de mi email'' equivale a decir qué email se
  está consultando --- \textbf{se pierde el anonimato} que se buscaba (el servidor sabe
  qué registro se está consultando, aunque no vea el email en texto plano).
  \item Si en cambio se usa \textbf{un único salt global} para toda la base, no cambia
  nada respecto de no usar salt: todos los emails se perturban con el mismo valor, así
  que el cliente puede seguir calculando el hash localmente sin pedir nada al servidor
  (el salt global es efectivamente parte de la constante pública del hash).
\end{itemize}
En síntesis: un salt por-cuenta \textbf{rompe el protocolo} (el cliente no puede
recalcular el hash sin filtrar información al servidor); un salt global no aporta nada
al anonimato porque equivale a no usar salt. En ningún caso mejora el anonimato deseado.

\textbf{b)} El cliente calcula $\mathrm{SHA256}(\text{email})$ \textbf{localmente} y
envía solo los primeros $x$ bits de ese hash. Por \textbf{resistencia a preimagen} de
SHA-256, el servidor no puede invertir esos $x$ bits para recuperar el email (ni
siquiera el hash completo, del cual solo recibe un prefijo). Además, como el servidor
responde con \textbf{todos} los hashes completos de la base que comparten ese prefijo
(no solo el que interesa), no puede saber cuál de esos hashes de la respuesta es el que
realmente le importa al cliente: la comparación final (¿aparece mi hash completo en la
lista recibida?) la hace el cliente \textbf{del lado suyo}, sin informarle nada al
servidor sobre el resultado. Por eso el servidor no puede saber ni qué email se
consultó (solo ve un prefijo compartido por muchos emails posibles) ni si la búsqueda
fue exitosa (esa comparación ocurre localmente en el cliente, nunca se envía de vuelta).

\textbf{c)} El enunciado define $k$-annonymity como la imposibilidad de asociar un dato
a una entidad discernible en un grupo de \textbf{al menos $k$ entidades}. Aplicado a
este esquema, se necesita que el prefijo de $x$ bits enviado por el cliente sea
compartido, en promedio, por al menos $k=256$ hashes de la base (para que el servidor
no pueda reducir el conjunto de sospechosos a un grupo menor a 256 emails). Si $N$ es la
cantidad de registros (emails) en la base y se asume que el hash distribuye
uniformemente sus salidas entre los $2^x$ prefijos posibles (consecuencia del efecto
avalancha del hash), la cantidad esperada de hashes por cada prefijo de $x$ bits es
$N/2^x$. El criterio para garantizar 256-anonimity es entonces pedir que esa cantidad
esperada sea al menos 256:
\begin{equation*}
\frac{N}{2^{x}} \ge 256 \quad\Longleftrightarrow\quad x \le \log_2\!\left(\frac{N}{256}\right).
\end{equation*}
Es decir, el valor $x$ (bits de prefijo a enviar) debe fijarse como
$\boxed{\,x \le \log_2(N/256)\,}$, con $N$ el tamaño de la base de datos: cuantos menos
bits de prefijo se envíen, más grande es el grupo de hashes indistinguibles entre sí (y
al revés, si $x$ crece, el grupo de hashes que matchean se achica y se pierde
anonimato). \textbf{Nota de fidelidad:} esta fórmula se deriva aplicando la definición
de $k$-anonimato dada en el propio enunciado junto con las propiedades de hash de la
materia (salida fija, efecto avalancha $\Rightarrow$ distribución aproximadamente
uniforme de prefijos); el concepto de $k$-anonimato en sí no figura en los resúmenes de
la materia, así que no se cita ningún resultado formal de la cátedra específico de
$k$-anonimato.

% =====================================================================
\subsection{Ejercicio 2 --- Red de sensores de oleoducto: CHACHA20 + HMAC-SHA3}
% =====================================================================

\begin{consigna}
Un oleoducto es una tubería que transporta petróleo a grandes distancias (a veces miles
de kilómetros). Dada la longitud, la dificultad de acceso en segmentos soterrados y el
costo de no poder utilizarlo, el diseño incluye un sistema de control activo que
permite identificar problemas en el transporte con una localización precisa de dónde
ocurre el problema.

Cada 200m se emplazan un conjunto de sensores dentro del oleoducto que monitorean la
presión y algunas características químicas que permiten estimar problemas como una fuga
o contaminación. Estos sensores están conectados físicamente con un cableado interno, y
cada 10 km hay concentradores, que son dispositivos externos a la tubería, que están
conectados tanto al segmento anterior como al posterior al dispositivo y recopilan la
información. Los concentradores transmiten por radiofrecuencia los resultados a una
central, donde son analizados para generar alarmas cuando sea necesario (por ejemplo,
si falta información de algún sensor).

Dado el valor del material transportado, un escenario de ataque contemplado incluye a
un grupo con acceso a tecnología sofisticada, capaz de tomar control físicamente de un
concentrador y manipularlo, con el objetivo de ``pinchar'' un tubo y extraer petróleo
en pequeñas cantidades (algunos camiones por día) para comercializarlo por canales no
oficiales.

Cada grupo de sensores tiene un identificador único público de 128 bits, un
identificador único privado de 128 bits (que puede ser utilizado, por ejemplo, como
clave, pero nunca es transmitido). Los sensores siguen un modelo de transmisión simple
estilo token-ring para hacer uso del cable que los conecta, y cuando es su turno,
transmiten un par (identificador, enc(valor sensado)), utilizando CHACHA20 como función
de cifrado de flujo.

Los concentradores cuentan con su propio par de identificadores, pero solo la central
conoce los identificadores privados del resto. Cada concentrador captura los mensajes
transmitidos por cada sensor durante 1 minuto, tanto de la sección anterior como la
posterior al dispositivo, y las envía en un único mensaje de la forma (identificador,
msg1|msg2|\ldots|msgn, mac(msg1|msg2|\ldots|msgn)), utilizando para el mac su
identificador privado y la función hmac-sha3.

Nota: a$|$b significa a seguido de b.

\begin{enumerate}
  \item[a)] Explicar qué debe hacer la central para verificar la integridad de los
  valores de cada sensor que recibió, antes de determinar si sus valores son normales o
  no. (1 pto)
  \item[b)] ¿Puede alguien con capacidad de reemplazar un concentrador (incluso
  averiguando su identificador privado) falsificar valores de sensores que la central
  de como validos? Explicar por qué no o como puede hacerlo. (1 pto)
  \item[c)] ¿Puede alguien con capacidad de reemplazar dos concentradores falsificar
  valores de sensores que la central de como validos? Explicar por qué no o como puede
  hacerlo. (1 pto)
\end{enumerate}
\end{consigna}

\paragraph{Temas.}
Primera Parte, Clase 2 --- Cifrado (\S Cifrado de flujo; \S Generador pseudoaleatorio;
\S no reutilizar la clave/keystream; \S los cifrados de flujo no son seguros bajo
múltiples cifrados con la misma clave). Primera Parte, Clase 3 --- MACs y cifrado
autenticado (\S Qué es un MAC; \S HMAC: MAC a partir de una función de hash, incluye
explícitamente HMAC-SHA3; \S infalsificabilidad / Mac-forge). Segunda Parte, Clase 8
--- Principios de diseño (\S mínimo privilegio / confinamiento, para el razonamiento
del inciso c).

\paragraph{Qué necesitás saber.}
\begin{itemize}
  \item \textbf{MAC (Message Authentication Code):} terna $(\mathrm{Gen},
  \mathrm{Mac}, \mathrm{Vrfy})$ con clave secreta compartida; $\mathrm{Vrfy}_k(m,t)=1$
  si y solo si $t$ es una etiqueta válida de $m$ bajo $k$. Sin la clave, no se puede
  generar ni verificar un tag válido (infalsificabilidad).
  \item \textbf{HMAC:} $\mathrm{HMAC}_k(m) = H\big((k\oplus\texttt{opad})\,\|\,
  H((k\oplus\texttt{ipad})\,\|\,m)\big)$; si $H$ es libre de colisiones, HMAC es
  infalsificable. La materia menciona explícitamente HMAC-SHA1, HMAC-MD5 (evitar) y
  \textbf{HMAC-SHA3}.
  \item \textbf{Cifrado de flujo (CHACHA20 es una instancia de esta familia, aunque el
  nombre concreto ``ChaCha20'' no aparece nombrado en los resúmenes de la materia):}
  $e_k(m) = G(k)\oplus m$, con $G$ un generador pseudoaleatorio. Es
  \textbf{determinístico}: la misma clave/semilla produce siempre el mismo keystream.
  \item \textbf{Los cifrados de flujo NO son seguros bajo múltiples cifrados con la
  misma clave} (como el OTP reusado): reusar la clave permite XOR-ear dos cifrados y
  cancelar el keystream, exponiendo relaciones entre los textos planos.
\end{itemize}

\paragraph{Notas y conexiones.}
El identificador privado del sensor actúa como \textbf{clave simétrica de cifrado de
flujo} para los datos, y el identificador privado del concentrador actúa como
\textbf{clave del HMAC} sobre el mensaje agregado. Son dos primitivas independientes
con dos claves independientes (una por sensor, para confidencialidad; una por
concentrador, para integridad/autenticación del agregado), lo que limita el
\textbf{alcance del compromiso} de cada clave --- la misma idea de compartimentar
privilegios que aparece en la unidad de Principios de diseño (mínimo privilegio,
confinamiento): comprometer una clave no debería dar poder sobre datos protegidos por
otra clave distinta.

\paragraph{Resolución.}

\textbf{a)} Para verificar la integridad del mensaje agregado recibido de un
concentrador, la central debe:
\begin{enumerate}
  \item Identificar el concentrador emisor por su \textbf{identificador público}
  (que sí viaja en el mensaje).
  \item Buscar el \textbf{identificador privado} de ese concentrador, que la central
  conoce (según el enunciado, ``solo la central conoce los identificadores privados del
  resto'').
  \item Recalcular ella misma $\mathrm{HMAC}_{k_{\text{priv}}}(\text{msg}_1\|\cdots\|
  \text{msg}_n)$ sobre la concatenación de mensajes recibida, usando el identificador
  privado del concentrador como clave y la función hmac-sha3 (misma construcción que
  usó el concentrador para generar el tag).
  \item Comparar el valor recalculado con el $\mathrm{mac}(\cdot)$ recibido en el
  mensaje: si coinciden, el mensaje \textbf{no fue alterado} en tránsito (íntegro y
  autenticado como proveniente de ese concentrador); si no coinciden, se descarta/marca
  como no confiable.
\end{enumerate}
Solo después de esa verificación (recalcular y comparar el HMAC) tiene sentido que la
central analice si los valores sensados son ``normales o no''; verificar antes evita
procesar datos manipulados como si fueran legítimos.

\textbf{b)} \textbf{Sí puede falsificar valores}, precisamente porque conoce el
identificador privado del concentrador reemplazado: ese identificador privado es
exactamente la \textbf{clave del HMAC} que la central usa para validar la integridad
del mensaje agregado de ese concentrador. Un atacante que reemplaza físicamente el
concentrador puede:
\begin{enumerate}
  \item Inventar o alterar los valores sensados de esa sección ($\text{msg}_1', \ldots,
  \text{msg}_n'$, p.ej.\ para esconder la caída de presión que delataría el
  ``pinchazo'').
  \item Calcular él mismo $\mathrm{HMAC}_{k_{\text{priv}}}(\text{msg}_1'\|\cdots\|
  \text{msg}_n')$ con el identificador privado que conoce (o ya lo tiene disponible por
  haber comprometido físicamente el dispositivo, según el enunciado).
  \item Enviar (identificador público, mensajes falsificados, HMAC recalculado con la
  clave correcta) a la central, que al recalcular el HMAC obtiene una coincidencia y da
  los valores por \textbf{válidos}.
\end{enumerate}
Es decir, el HMAC protege contra un atacante que \textbf{no} conoce la clave (garantiza
infalsificabilidad frente a un adversario sin oráculo/clave), pero no protege si el
propio concentrador --- y por lo tanto quien lo controla --- es el poseedor legítimo de
esa clave: quien tiene la clave puede generar tags válidos para cualquier mensaje que
elija, exactamente como el emisor legítimo.

\textbf{c)} Comprometer \textbf{dos} concentradores le da al atacante control sobre dos
identificadores privados (dos claves de HMAC) en vez de una, pero cada concentrador
usa \textbf{su propio par de identificadores}, independiente del resto (según el
enunciado, ``los concentradores cuentan con su propio par de identificadores''). Por lo
tanto:
\begin{itemize}
  \item El atacante puede falsificar, de la misma manera que en b), los valores de
  \textbf{los dos segmentos} cubiertos por esos dos concentradores comprometidos (el
  doble de alcance que en b), pero nada más).
  \item \textbf{No puede} falsificar valores de sensores/concentradores fuera de esas
  dos secciones: no conoce los identificadores privados del resto de los
  concentradores ni de los sensores de otras secciones, y sin esa clave no puede
  generar un HMAC que la central acepte como válido para esos otros segmentos (por la
  infalsificabilidad del HMAC bajo una clave que el atacante desconoce).
\end{itemize}
En síntesis: comprometer $m$ concentradores solo compromete la integridad de esos $m$
segmentos, nunca la del resto de la red --- consecuencia directa de que cada
concentrador (y cada sensor) tiene su propia clave independiente, limitando el
``radio de explosión'' de cada compromiso individual (mismo principio de aislar
privilegios/claves visto en la unidad de Principios de diseño).

% =====================================================================
\subsection{Ejercicio 3 --- Hipótesis de falla para pentest de app móvil de exchange}
% =====================================================================

\begin{consigna}
Una exchange de criptomonedas está expandiendo su aplicación web actual con una nueva
aplicación móvil nativa para Android y otra para iOS.

Se sabe que:
\begin{itemize}
  \item La aplicación se comunica con el servidor mediante un API REST.
  \item La aplicación requiere que el usuario se registre con un email y una
  contraseña.
  \item Sin autenticarse en la aplicación, se puede ver un resumen del valor actual de
  compra y venta de cada criptomoneda listada, y datos de volumen de operaciones.
  \item Para ingresar dinero se puede utilizar una tarjeta de crédito o una cuenta
  bancaria. Es necesario registrar los datos del medio a utilizar. Para retirar dinero,
  solo se puede utilizar una transferencia bancaria.
  \item Los ingresos y salidas de dinero están gestionados por un sistema de pagos
  externo, AlphaPay, que es quien procesa las operaciones. El proveedor disponibiliza un
  api para llevar a cabo las transacciones. Dicha API procesa las transacciones, pero no
  almacena información de medios de pago, que sí almacena la plataforma para no
  requerir al usuario ingresar todos los datos de cuenta en cada transacción.
  \item La aplicación notifica por email cuando una orden de compra o venta se ejecuta,
  incluyendo el resumen de la misma.
\end{itemize}

Ante un pedido de realizar un pentest que abarque la app, establecer 4 hipótesis de
falla \textbf{pertinentes} para el caso, basadas en problemas de seguridad en
aplicaciones. Explicar \textbf{para cada una}, como probaría según la metodología (la
acción concreta que ejecutaría para verificar o refutar la hipótesis). (4 puntos, 1 pto
cada uno)

\textbf{Aclaración:} No utilizar problemas genéricos como hipótesis de falla. Plantear
situaciones y lugares concretos donde podría estar el problema. Por ejemplo, en lugar de
decir ``puede haber un buffer overflow'', decir ``es probable que haya un buffer
overflow en la lectura del parámetro x que debe interpretarse como número y puede tener
un buffer pequeño asumiendo que el número tiene menos de los dígitos de un long''.
\end{consigna}

\paragraph{Temas.}
Segunda Parte, Clase 6 --- Pentesting (\S Metodología de Hipótesis de Falla, pasos y
ejemplos de hipótesis concretas: seguridad en APIs, relaciones de confianza). Segunda
Parte, Clase 3 --- Principios de diseño y vulnerabilidades (\S mediación completa;
\S límite de confianza / trust boundary; \S CWE Top 25: injection, buffer overflow).
Segunda Parte, Clase 2 --- Autenticación (\S fuerza bruta / throttling). Segunda Parte,
Clase 7 --- Protección de datos (\S PCI-DSS: acotar datos críticos, tokenización).

\paragraph{Qué necesitás saber.}
\begin{itemize}
  \item La \textbf{metodología de hipótesis de falla} exige, para cada hipótesis,
  describir cómo se \textbf{probaría} (verificaría o refutaría): analizando
  documentación/comportamiento, o --- si es necesario --- intentando explotar la
  vulnerabilidad de forma controlada y repetible/concluyente.
  \item \textbf{Mediación completa / trust boundary:} todo dato que cruza un límite de
  confianza (ej.\ del cliente --- app móvil, que el atacante controla totalmente --- al
  servidor) debe validarse en el lado de \emph{mayor} confianza; no se puede confiar en
  valores calculados o enviados por el cliente.
  \item \textbf{PCI-DSS (protección de datos):} acotar la cantidad de datos críticos
  almacenados según su criticidad (no proteger todo por igual), y usar
  \textbf{tokenización} para que, si se filtran los datos guardados por la plataforma,
  no tengan valor por sí solos.
  \item \textbf{Fuerza bruta / throttling:} sin límite de intentos o política de
  contraseñas, un endpoint de autenticación es vulnerable a ataques de fuerza
  bruta/diccionario contra cuentas existentes.
\end{itemize}

\paragraph{Notas y conexiones.}
El enunciado da varias pistas textuales que apuntan directamente a puntos débiles ya
vistos en la materia: (1) la plataforma \textbf{almacena datos de medios de pago}
(tarjeta/cuenta bancaria) para no pedirlos en cada transacción --- esto es justo lo que
PCI-DSS busca acotar/tokenizar, y es un dato crítico guardado localmente en vez de
delegado al proveedor externo (AlphaPay), que explícitamente \emph{no} los almacena;
(2) el flujo de compra/venta involucra un valor de cotización y un monto que el cliente
podría manipular si el backend no revalida contra el servidor/AlphaPay (mediación
completa); (3) el registro con email/contraseña es un endpoint de autenticación clásico
para fuerza bruta si no hay throttling; (4) los emails de notificación de operaciones
incluyen un ``resumen'' cuyo contenido podría filtrar de más.

\paragraph{Resolución.}

\textbf{Hipótesis 1 --- Almacenamiento inseguro de datos de medios de pago (PCI-DSS).}
Es probable que los datos del medio de pago (número de tarjeta de crédito o datos de la
cuenta bancaria) que la plataforma guarda para no pedirlos en cada transacción se
almacenen en la base de datos de la aplicación sin la protección exigida por PCI-DSS
(sin tokenización, es decir, guardando el dato real en vez de un token que apunte a él
en un sistema aparte ultra-protegido) y/o asociados incorrectamente al usuario,
permitiendo que un medio de pago registrado por un usuario sea reutilizable desde otra
cuenta.
\emph{Prueba:} con dos cuentas de prueba, registrar un medio de pago en la cuenta A y
luego, desde la cuenta B, intentar usar el identificador interno de ese medio de pago
(obtenido por ejemplo interceptando el tráfico de la API REST de la cuenta A) en una
operación de ingreso de dinero; verificar si la transacción se acredita/ejecuta contra
el medio de pago de otro usuario. Adicionalmente, revisar (con acceso de caja
gris/blanca) si el campo del medio de pago se guarda en texto plano o tokenizado en la
base.

\textbf{Hipótesis 2 --- Confianza indebida en valores enviados por el cliente al
ejecutar una orden.} Es probable que, al confirmar una compra o venta de una
criptomoneda, el backend confíe en el monto y/o el precio de cotización enviados desde
la app móvil por el API REST al iniciar la orden, en lugar de revalidar esos valores
contra la cotización real del servidor en el momento de ejecutar la operación --- el
cliente móvil cruza el límite de confianza (trust boundary) hacia el servidor y ese es
el punto donde debería mediarse completamente la validación.
\emph{Prueba:} interceptar (con un proxy) la request que envía la app al confirmar una
orden de compra/venta, modificar el parámetro de precio o monto antes de que llegue al
servidor, y verificar si la orden se ejecuta con el valor alterado en vez de con la
cotización real vigente en el momento de la ejecución.

\textbf{Hipótesis 3 --- Falta de límite de intentos (rate limiting) en el endpoint de
registro/login con email y contraseña.} Es probable que el endpoint de autenticación
(registro/login con email y contraseña) no implemente throttling ni bloqueo de cuenta
tras varios intentos fallidos, permitiendo un ataque de fuerza bruta o de diccionario
contra contraseñas de cuentas existentes.
\emph{Prueba:} realizar múltiples intentos de autenticación consecutivos contra una
cuenta de prueba válida usando contraseñas incorrectas (diccionario de contraseñas
comunes) y observar si el sistema bloquea, limita (con backoff) o registra/alerta sobre
los intentos, o si permite continuar indefinidamente.

\textbf{Hipótesis 4 --- Exceso de información sensible en el email de notificación de
operaciones.} Es probable que el email que notifica la ejecución de una orden de compra
o venta, al incluir ``el resumen de la misma'', exponga información sensible en exceso
(por ejemplo datos parciales del medio de pago usado, montos exactos de fondos
disponibles, u otros identificadores internos de la cuenta/orden), quedando dicha
información expuesta en un canal (el email) menos controlado que el API REST de la
aplicación.
\emph{Prueba:} ejecutar una orden de compra/venta de prueba y analizar el contenido
completo del email recibido, verificando si incluye datos sensibles (número parcial o
completo de tarjeta/cuenta, identificadores internos, saldo total) que excedan lo
estrictamente necesario para informar la operación.

\medskip
\noindent\textbf{Nota de fidelidad:} las cuatro hipótesis se construyen a partir de
datos concretos del propio enunciado (qué se registra, qué almacena la plataforma, qué
notifica) combinados con los conceptos de la materia citados en ``Qué necesitás saber''
(PCI-DSS/tokenización, mediación completa/trust boundary, throttling). No se incluyen
hipótesis genéricas de OWASP Top Ten que no estén ancladas a un dato concreto del
enunciado, siguiendo la aclaración del examen.

\clearpage
%=====================================================================
\section{Final --- 10/07/2023}
%=====================================================================

%=====================================================================
\subsection{Ejercicio 1 --- Sensor $\to$ maestro con AES-CBC: integridad y replay}
%=====================================================================

\begin{consigna}
Una empresa tiene un sistema donde hay varios sensores que comunican a un dispositivo
maestro sus diferentes niveles de sensado. Intercambian claves con Diffie-Hellman en su
primera comunicación, y a partir de eso en todos los mensajes que los sensores envían
(encriptados con AES-CBC) tienen la siguiente estructura:
\begin{equation*}
\#\text{sensor} \,\|\, e_k(\#\text{sensor} \,\|\, \text{valor\_sensado} \,\|\, \text{datos\_adicionales})
\end{equation*}
\begin{enumerate}
  \item[1.] Hay un atacante que hace de alguna forma que los mensajes se alteren y
  tengan datos adicionales no válidos, generando que el dispositivo maestro tire una
  señal de alerta. \textquestiondown{}Cómo es posible? Encontrar una manera de solucionar este error en
  el protocolo sin nuevas claves. (1,5 puntos)
  \item[2.] Hay otro atacante que está enviando repetidamente la señal de sensado de un
  sensor, haciendo que el dispositivo maestro registre niveles de sensado que no son
  correctos. \textquestiondown{}Cómo es posible? Encontrar una forma de modificar el mensaje para que
  esto no ocurra sin usar nuevas claves. (1,5 puntos)
\end{enumerate}
\end{consigna}

\paragraph{Temas.}
Primera Parte, Clase de MACs y Cifrado Autenticado (\S El nuevo escenario de ataque:
integridad; \S Maleabilidad; \S Cifrado autenticado: Encrypt-then-MAC); Clase de
Protocolos (\S Nonce y frescura; \S Replay; \S Diffie-Hellman como intercambio de clave).

\paragraph{Qué necesitás saber.}
\begin{itemize}
  \item \textbf{El cifrado por sí solo no da integridad.} El material lo remarca con el
  ejemplo del sistema de RR.HH.: un criptosistema CPA-Secure (incluso uno ``bueno'') no
  impide que un atacante \emph{manipule} el texto cifrado sin conocer la clave. La cita
  del profesor: \prof{En general el problema de integridad no es evitar la modificación,
  porque muchas veces no se puede, sino detectarla.}
  \item \textbf{Maleabilidad:} propiedad de que modificar el texto cifrado se traduce de
  forma sistemática/controlable en una modificación del texto plano. En un cifrado de
  flujo, $c'=c\oplus x \Rightarrow \mathrm{Dec}_k(c')=m\oplus x$. AES-CBC también es
  maleable: modificar un bloque de cifrado altera de forma predecible (vía el XOR con el
  bloque anterior) el bloque de texto plano correspondiente, y además \textbf{corrompe}
  el bloque siguiente al descifrar --- pero eso no evita que el atacante logre inyectar
  datos con sentido si conoce o controla estructura del mensaje.
  \item \textbf{Ningún criptosistema visto es CCA-Secure} (recuadro ``destacado'' del
  material): el cifrado está pensado para confidencialidad, no para resistir
  manipulación. Por eso la integridad requiere \textbf{otra primitiva: el MAC}.
  \item \textbf{Solución sin cambiar claves --- Cifrado autenticado / MAC:} el material
  define la construcción \textbf{Encrypt-then-MAC} como la única con demostración
  general de seguridad: se cifra $m$ obteniendo $c$, y se etiqueta con
  $t=\mathrm{Mac}_{k_2}(c)$ con clave \textbf{independiente} $k_2\ne k_1$; el receptor
  \textbf{primero verifica} el tag y solo si es válido descifra. Alternativamente (y sin
  agregar una clave nueva, ya que acá solo se cuenta con $k$ derivada de DH) se puede usar
  un \textbf{HMAC con una clave derivada de la misma $k$} (p.\ ej.\ vía un KDF/hash, tal
  como aparece en el ejercicio análogo del material: $c=E_{k_1}(m\|H(k_2\|m))$, donde
  $k_1,k_2$ salen de la misma clave maestra) para no necesitar un nuevo intercambio DH.
  \item \textbf{Nonce, timestamp y replay (Clase de Protocolos):} un \textbf{nonce}
  (number used once) es un número aleatorio de un solo uso que aporta frescura y permite
  detectar/evitar \textbf{replays} (repetición de mensajes viejos); el otro mecanismo de
  frescura es el \textbf{timestamp}. El material lo usa en Needham-Schroeder y en la
  variante Denning-Sacco justamente para evitar el reenvío de mensajes/claves viejas.
\end{itemize}

\paragraph{Notas y conexiones.}
El enunciado describe un mensaje cifrado con AES-CBC pero \textbf{sin ningún mecanismo
de integridad ni de frescura}: no hay MAC sobre el mensaje, y no hay nonce, contador ni
timestamp dentro de lo cifrado. Esto es exactamente el escenario que la unidad de MACs
identifica como el ``nuevo escenario de ataque'': un cifrado CPA-Secure (o incluso
CCA-inseguro, que es la realidad de todos los criptosistemas vistos) puede ser
manipulado sin que el atacante conozca la clave. La restricción ``sin nuevas claves''
en ambos incisos apunta a agregar mecanismos (MAC derivado de la misma clave, contador)
en vez de rehacer el intercambio Diffie-Hellman.

\paragraph{Resolución.}

\textbf{1) Por qué es posible y cómo solucionarlo.}
AES-CBC \textbf{cifra pero no autentica}: es un modo de cifrado por bloques que da
confidencialidad (si el criptosistema es CPA-Secure), pero --- como remarca el
material --- \textbf{ningún} criptosistema visto es CCA-Secure, es decir, ninguno
resiste manipulación del texto cifrado. Un atacante que intercepta el mensaje puede
alterar directamente el texto cifrado $e_k(\cdots)$ (por ejemplo, invirtiendo bits de
algún bloque cifrado, lo cual --- por la estructura de CBC --- modifica de forma
predecible/controlable el bloque de texto plano correspondiente al descifrar, aunque
corrompa el bloque siguiente). Como no hay ningún chequeo de integridad, el dispositivo
maestro descifra ``basura'' que cae en el campo \texttt{datos\_adicionales} sin que el
receptor pueda detectar que fue modificado --- eso es lo que dispara la alerta (datos
adicionales no válidos), sin que el atacante haya necesitado la clave.

\textbf{Solución (sin nuevas claves):} agregar integridad con un \textbf{MAC/HMAC}
calculado sobre el mensaje (siguiendo el patrón Encrypt-then-MAC del material, que es el
único con demostración de seguridad), usando una clave derivada de la misma $k$ obtenida
por Diffie-Hellman (no hace falta un nuevo intercambio):
\begin{equation*}
c = e_k(\#\text{sensor}\,\|\,\text{valor\_sensado}\,\|\,\text{datos\_adicionales}) \,\|\, m,
\qquad m = \mathrm{HMAC}_k(c)
\end{equation*}
El maestro, al recibir $(\#\text{sensor}, c, m)$, \textbf{primero verifica} $m$ contra
$c$ recalculando el HMAC; si no coincide, descarta el mensaje directamente (nunca llega
a descifrar datos manipulados). Solo si el MAC es válido, descifra y usa
\texttt{datos\_adicionales}. Así cualquier alteración del cifrado es detectada antes de
generar una alerta espuria.

\textbf{2) Por qué es posible y cómo solucionarlo.}
El segundo atacante no modifica el contenido: \textbf{reenvía (replay)} un mensaje
legítimo capturado anteriormente. Como la estructura del mensaje no incluye ningún
elemento de \textbf{frescura} (ni nonce, ni contador, ni timestamp), el dispositivo
maestro no tiene forma de distinguir un mensaje nuevo de uno viejo reenviado, y aun con
un MAC válido (agregado en el punto 1) el mensaje \emph{reenviado} seguiría verificando
correctamente --- el MAC da integridad y autenticación de origen, pero no dice nada
sobre si el mensaje es \emph{nuevo}. Por eso el maestro termina registrando niveles de
sensado viejos como si fueran actuales.

\textbf{Solución (sin nuevas claves):} agregar un \textbf{contador (o timestamp)} dentro
del contenido protegido por el MAC, tal como el material propone con la variante
Denning-Sacco (agregar un timestamp $T$ dentro del bloque protegido para impedir
reutilizar un mensaje/ticket viejo):
\begin{equation*}
c = e_k(\#\text{sensor}\,\|\,\text{valor\_sensado}\,\|\,\text{datos\_adicionales}\,\|\,n_i)\,\|\,m,
\qquad m=\mathrm{HMAC}_k(c)
\end{equation*}
donde $n_i$ es un contador estrictamente creciente por sensor (o un timestamp). El
maestro guarda el último contador válido recibido de cada sensor y \textbf{rechaza}
cualquier mensaje cuyo contador/timestamp no sea mayor al último aceptado, aun cuando el
MAC verifique correctamente. Esto detecta y descarta el replay sin necesidad de repetir
el intercambio Diffie-Hellman ni generar nuevas claves.

%=====================================================================
\subsection{Ejercicio 2 --- Alternativas de MAC para datos grandes (árbol de Merkle)}
%=====================================================================

\begin{consigna}
Se quiere enviar un conjunto de datos muy grande. Si bien se puede mandar la
información junto con un MAC de la información entera, en caso de que una parte del
mensaje tenga un error, es necesario mandar el mensaje entero de vuelta. Se buscan
alternativas para MAC, donde se parte la data en bloques más pequeños, de forma que
permitan no tener que enviar la información entera en caso de una falla en uno de los
bloques. Las alternativas son las siguientes:
\begin{enumerate}
  \item[1.] En vez de calcular el MAC de cada bloque, se calcula el MAC' que es: el
  contenido del bloque anterior concatenado al MAC del bloque actual. Excepto en el
  bloque cero, que como no tiene anterior, se ponen 0s. (1 punto)
  \item[2.] Se calcula el MAC de la concatenación de los MAC de cada bloque. (1 punto)
  \item[3.] Se calcula el MAC usando el árbol de Merkle (buscar ejemplo de árbol de
  Merkle). (1 punto)
\end{enumerate}
Para cada caso evaluar si tiene integridad y en caso de que sí, cómo harían cuando un
bloque se envía con un dato erróneo (si es que hay una alternativa a no enviar todo el
mensaje de vuelta).
\end{consigna}

\paragraph{Temas.}
Primera Parte, Clase de MACs y Cifrado Autenticado (\S Qué es un MAC; \S Falsificación
por bloques (parcial 2015); \S CBC-MAC; \S HMAC; \S Modelo iterativo Merkle-Damgård).

\begin{warnbox}{Aviso de fidelidad --- el árbol de Merkle NO está en el material de la materia}
El enunciado dice explícitamente ``buscar ejemplo de árbol de Merkle'': el \textbf{árbol
de Merkle (hash tree)} como estructura de datos \textbf{no aparece} en los resúmenes de
la materia (ni en la unidad de MACs ni en ninguna otra). Lo que sí está en el material es
la construcción \textbf{Merkle-Damgård} (el esquema iterativo genérico con el que se
construyen las funciones de hash como SHA-2, encadenando $H_i=f(H_{i-1},x_i)$) --- que es
un concepto \textbf{distinto} (sirve para construir una función de hash sobre un mensaje
secuencial, no es una estructura de árbol para autenticar bloques independientes). Por lo
tanto, la resolución de la alternativa 3 usa la noción \emph{general} de MAC (Clase de
MACs: $\mathrm{Mac}_k$ sobre una función de hash / HMAC) aplicada de forma recursiva sobre
pares de hojas, que es la idea estándar de un árbol de hashes, pero \textbf{esa
estructura en árbol específica no proviene de las clases provistas}: se avisa siguiendo
la regla de oro en vez de presentarla como si estuviera dictada.
\end{warnbox}

\paragraph{Qué necesitás saber.}
\begin{itemize}
  \item \textbf{Qué exige un MAC (Clase de MACs):} terna $(\mathrm{Gen},\mathrm{Mac},
  \mathrm{Vrfy})$; $\mathrm{Vrfy}_k(m,t)=1$ si $t$ es la etiqueta correcta de $m$. Un MAC
  \textbf{robusto siempre procesa la totalidad} del contenido que dice proteger: si deja
  algo afuera, ese resto queda \textbf{sin proteger} y es falsificable modificándolo
  libremente.
  \item \textbf{Falsificación por bloques (parcial 2015, en el material):} cuando un MAC
  se define como bloques \textbf{calculados de forma independiente}
  ($\mathrm{Mac}_k(m)=\langle F_k(m_1), F_k(F_k(m_2))\rangle$ y similares), el ataque
  típico es \textbf{reusar/intercambiar bloques} entre mensajes distintos para forjar un
  tag válido de otro mensaje --- justamente \emph{porque} cada bloque se verifica de
  forma independiente. Esto es una debilidad de \emph{seguridad} (permite forjar
  combinando piezas de mensajes distintos con MACs válidos), pero la misma independencia
  es la que permite, del lado de la \textbf{detección de errores de transmisión}
  (no forjas activas, sino bloques corrompidos), \textbf{verificar cada bloque por
  separado} y así saber exactamente cuál falló.
  \item \textbf{CBC-MAC (Clase de MACs):} $t_i=F_k(t_{i-1}\oplus m_i)$, el tag final
  ``depende'' de todos los bloques anteriores por el encadenamiento, pero solo se
  conserva/verifica el \textbf{último} tag $t_j$: no hay forma de saber, a partir de un
  único CBC-MAC roto, cuál bloque intermedio fue el que falló (hay que volver a mandar
  todo el mensaje para poder recalcular y comparar).
\end{itemize}

\paragraph{Notas y conexiones.}
El punto clave a verificar es la alternativa 2. En el material, la idea de que un MAC
\textbf{calculado independientemente por bloque} permita identificar exactamente qué
bloque fue alterado es consistente con la sección de ``Falsificación por bloques'': ahí
el material señala que esos esquemas son atacables \emph{justamente porque} cada
$F_k(m_i)$ se computa y se verifica de forma aislada del resto. Esa misma propiedad
(verificación aislada por bloque) es la que, en el escenario de este ejercicio
(corrupción/error de un bloque, no un atacante activo forjando), permite decir
\textbf{cuál bloque en particular no verifica}. Por lo tanto, el razonamiento de que
"si falla un bloque hay que reenviar todo el mensaje porque no se puede saber cuál
bloque falló" para la alternativa 2 \textbf{no es correcto} y se corrige en la
resolución de abajo.

\paragraph{Resolución.}

\textbf{1) MAC$'_i$ = contenido del bloque anterior $\|$ MAC del bloque actual
(MAC$'_i = B_{i-1}\,\|\,\mathrm{Mac}(B_i)$, con $B_{-1}=$ ceros).}
\begin{itemize}
  \item \textbf{\textquestiondown{}Tiene integridad?} Sí, para el contenido de cada bloque: cada
  $\mathrm{Mac}(B_i)$ se calcula (según la definición del material) sobre $B_i$ y se
  puede verificar de forma independiente por bloque, exactamente igual que en la
  alternativa 2. Además, al incluir el contenido crudo de $B_{i-1}$ junto al tag de
  $B_i$, el receptor puede \textbf{cotejar que el bloque anterior recibido coincida}
  con el que está referenciado en $\mathrm{MAC}'_i$: esto agrega una verificación de
  \textbf{orden/posición} (detecta si se reordenaron o eliminaron bloques) que la
  alternativa 2 no ofrece por sí sola, a costa de duplicar el contenido de cada bloque
  (se transmite $B_{i-1}$ dos veces).
  \item \textbf{Ante un bloque con error:} como $\mathrm{Mac}(B_i)$ se calcula solo
  sobre $B_i$ (no está encadenado criptográficamente con los bloques anteriores, a
  diferencia del CBC-MAC), el fallo se localiza en el bloque $i$ cuyo
  $\mathrm{Mac}(B_i)$ no verifica. Alcanza con \textbf{reenviar únicamente el bloque
  $B_i$} (el receptor ya tiene, del propio $\mathrm{MAC}'_{i+1}$, la copia correcta del
  contenido de $B_i$ para volver a verificar la cadena de orden una vez corregido).
\end{itemize}

\begin{warnbox}{Corrección respecto de una lectura errónea}
No corresponde afirmar que, si falla un bloque, ``hay que reenviar desde ese bloque y el
siguiente'': dado que $\mathrm{Mac}(B_i)$ en esta construcción se calcula únicamente
sobre $B_i$ (no sobre $B_i$ encadenado con $B_{i-1}$ dentro de la función de MAC), la
localización del error es exacta en el bloque $i$, y $\mathrm{MAC}'_{i+1}$ no cambia
al corregir $B_i$ (ya traía la copia correcta desde el envío original).
\end{warnbox}

\textbf{2) MAC = MAC(B$_1$) $\|$ $\cdots$ $\|$ MAC(B$_n$).}
\begin{itemize}
  \item \textbf{\textquestiondown{}Tiene integridad?} Sí: si se altera cualquier bloque $B_i$, su
  $\mathrm{Mac}(B_i)$ correspondiente deja de verificar (asumiendo un MAC infalsificable
  como los del material: CBC-MAC o HMAC), mientras que los tags de los demás bloques
  siguen siendo válidos.
  \item \textbf{Ante un bloque con error --- corrección del razonamiento del alumno:}
  \textbf{sí es posible saber exactamente cuál bloque falló}, y por lo tanto \textbf{no
  hace falta reenviar el mensaje entero}. Como cada $\mathrm{Mac}(B_i)$ se computa y se
  verifica de forma \textbf{independiente} (el mismo punto que el material remarca en
  ``Falsificación por bloques'': cada bloque se calcula por separado), el receptor
  recalcula $\mathrm{Mac}(B_i)$ para cada $i$ y lo compara contra la posición $i$ de la
  cadena recibida; el (o los) bloque(s) cuyo MAC no coincide son exactamente los que
  llegaron corruptos. Alcanza con \textbf{reenviar solo ese/esos bloque(s)}.
  \item Este es justo el punto que había que revisar contra el material: la
  independencia de los MAC por bloque, que en ``Falsificación por bloques'' se señala
  como una \textbf{debilidad de seguridad} (permite reusar/intercambiar bloques entre
  mensajes distintos para forjar un tag), es la \textbf{misma propiedad} que --- en este
  contexto de detección de errores de transmisión, sin un atacante activo forjando ---
  permite \textbf{localizar con precisión} el bloque dañado. No hay contradicción: son
  dos consecuencias de la misma independencia, una negativa (falsificabilidad por
  recombinación) y una positiva (localización de errores) en este ejercicio.
\end{itemize}

\textbf{3) MAC calculado con un árbol de Merkle.}
Como se indicó en el aviso de fidelidad de arriba, la estructura específica del árbol de
Merkle no está desarrollada en el material de la materia (solo Merkle-Damgård, que es
otra cosa); se resuelve con la idea general de hash/MAC recursivo sobre pares de bloques,
avisando que el detalle de la estructura en árbol excede las clases provistas:
\begin{itemize}
  \item Cada hoja se etiqueta con el hash/MAC de su bloque: $h_i = \mathrm{Mac}(B_i)$
  (o $H(B_i)$ si se usa una función de hash del material, p.\ ej.\ SHA-3). Cada nodo
  interno combina sus dos hijos: $h = \mathrm{Mac}_k(h_{\text{izq}}\|h_{\text{der}})$,
  hasta llegar a una única \textbf{raíz} que resume todo el conjunto de datos.
  \item \textbf{\textquestiondown{}Tiene integridad?} Sí: si se modifica cualquier bloque hoja, su hash
  cambia, y ese cambio se propaga hacia arriba modificando todos los hashes de los
  nodos ancestros hasta la raíz --- cualquier alteración es detectable comparando la
  raíz recalculada contra la raíz original.
  \item \textbf{Ante un bloque con error:} \textbf{no hace falta reenviar todo el
  mensaje.} Alcanza con reenviar el bloque hoja corregido junto con los hashes
  ``hermanos'' del camino desde esa hoja hasta la raíz (un camino de altura
  $\log_2 n$), de forma que el receptor pueda recalcular únicamente esa rama y
  verificar que la nueva raíz coincide con la original. Esto es más eficiente en ancho
  de banda de verificación que la alternativa 2 cuando $n$ es muy grande (overhead
  logarítmico en vez de lineal en la cantidad de bloques), aunque ambas permiten
  localizar y corregir un único bloque sin reenviar el mensaje completo.
\end{itemize}

\textbf{Resumen comparativo:} las tres alternativas dan integridad y, contrariamente a
enviar un único MAC sobre el mensaje entero (que obliga a reenviar todo ante cualquier
error), las tres permiten identificar el bloque afectado y reenviar solo lo necesario:
la 1 y la 2 con overhead lineal (un tag por bloque, más en la 1 el contenido duplicado
del bloque anterior para chequeo de orden), y la 3 (árbol de Merkle, fuera del material
provisto pero pedida explícitamente por el enunciado) con overhead logarítmico por
verificación.

%=====================================================================
\subsection{Ejercicio 3 --- Vulnerabilidades de un sitio de noticias}
%=====================================================================

\begin{consigna}
Se tiene una página de noticias donde hay usuarios que pueden ratear una noticia
positiva o negativamente. Los publicadores de noticias pueden subir noticias
pegándola a una API REST. Los usuarios interactúan con la página a través de una
página web que se comunica con la API. Los ingresos de las publicidades de las
noticias se pagan a los creadores según la cantidad de visitas que tiene una noticia.

Se encontraron 4 vulnerabilidades:
\begin{enumerate}
  \item Para evaluar la positividad de una noticia, los usuarios tocan un botón que
  hace un GET a \texttt{/\{news\_id\}/positive} o \texttt{/\{news\_id\}/positive/negative},
  pero se encontró que se puede hacer SSRF a través de esto.
  \item El contenido de una noticia se postea tal cual como se recibe, sin hacer
  validaciones.
  \item Las cookies de sesión no son \texttt{http\_only}, y se pueden acceder a través
  de javascript.
  \item Un publicador de noticia hace un llamado al backend con una API KEY que se
  introduce en el pedido de la siguiente forma: \texttt{/api/\{API\_KEY\}/publish}.
\end{enumerate}
\begin{enumerate}
  \item[a)] Explicar cómo podría un atacante sumarse dinero de forma maliciosa con los
  puntos 2 y 3. (1,5 puntos)
  \item[b)] Explicar cómo tomaría ventaja un atacante de la vulnerabilidad 1. (1,5 puntos)
  \item[c)] Explicar por qué es inseguro que se envíen las API KEYS como lo indica el
  punto 4 e indicar cómo lo corregirían. (1 punto)
\end{enumerate}
\end{consigna}

\paragraph{Temas.}
Segunda Parte, Clase de Principios de Diseño y Vulnerabilidades (\S Vulnerabilidades
web: XSS, CSRF, XXE; \S Buffer overflow e inyección de código, como marco general de
``mezclar datos con código/control'').

\begin{warnbox}{Aviso de fidelidad --- varios términos del enunciado no están (con ese nombre) en el material}
\begin{itemize}
  \item \textbf{SSRF (Server-Side Request Forgery)} \textbf{no aparece con ese nombre}
  en los resúmenes de la materia. Lo más cercano en el material es \textbf{XXE (XML
  External Entities)}: ``XML permite cargar esquemas desde una URL; si esa URL se usa
  para hacer un GET a otra página, un atacante puede escanear la red interna visible por
  el servidor''. Es el mismo patrón conceptual (el servidor hace, por instrucción del
  atacante, un request a un destino que no controla el usuario) pero el material lo
  desarrolla solo para el caso XXE, no como SSRF general vía un endpoint REST. Se marca
  para que quede claro que la resolución de b) se apoya en esa analogía y no en un
  contenido de SSRF dictado explícitamente.
  \item El atributo de cookie \textbf{\texttt{http\_only}} \textbf{no se menciona
  explícitamente} en el material (sí se cubre, en la unidad de Vulnerabilidades web, que
  XSS permite ``robar la sesión'' ejecutando JS en el browser de la víctima, que es la
  consecuencia directa de no tener esa protección).
  \item Los mecanismos concretos de \textbf{transporte seguro de API keys} (headers,
  OAuth/JWT) y el problema de que un secreto en la URL quede en \textbf{logs, historial
  del navegador o se filtre por el header \texttt{Referer}} \textbf{tampoco están
  desarrollados explícitamente} en el material provisto. La resolución de c) se apoya en
  los principios generales de diseño del material (mediación completa, fallar de forma
  segura, mínimo privilegio) más el criterio general de no exponer secretos donde quedan
  registrados, pero se avisa que el material no cubre el detalle técnico de dónde
  específicamente queda expuesta una URL.
\end{itemize}
\end{warnbox}

\paragraph{Qué necesitás saber.}
\begin{itemize}
  \item \textbf{XSS (CWE-79, Clase de Vulnerabilidades web):} inyección de JavaScript
  que termina \textbf{ejecutado por el browser de la víctima}, aprovechando la confianza
  entre el sitio y la sesión del usuario; ocurre cuando se hace render de HTML
  mezclando un dato controlado por el usuario con código sin escaparlo. Consecuencia
  remarcada en el material: \textbf{permite robar la sesión}. Defensa: sanitizar/escapar
  toda salida (\emph{escape on output}).
  \item \textbf{``El contenido se postea tal cual, sin validaciones''} (vulnerabilidad 2)
  es exactamente la falta de sanitización de entrada/salida que el material describe
  como causa de XSS: el dato del publicador (potencialmente conteniendo
  \texttt{<script>}) queda embebido en la página que ven los demás usuarios sin
  escapar --- un \textbf{XSS almacenado} (\emph{stored XSS}), ya que el payload queda
  guardado en la noticia y se ejecuta cada vez que otro usuario la visita.
  \item \textbf{XXE (Clase de Vulnerabilidades web):} el servidor termina haciendo un
  request (GET) a una URL/destino determinado por un input controlado por el atacante,
  lo que permite usar al servidor como intermediario para alcanzar recursos que el
  atacante no podría alcanzar directamente (p.\ ej.\ la red interna).
  \item \textbf{Principios de Saltzer \& Schroeder (Clase de Principios de Diseño):}
  \emph{mediación completa} (todo acceso debe pasar por un chequeo, sin atajos) y
  \emph{fallar de forma segura} son relevantes para pensar por qué exponer un secreto en
  una URL es delicado: la URL viaja/queda registrada por canales que no están pensados
  para llevar secretos.
\end{itemize}

\paragraph{Notas y conexiones.}
Las vulnerabilidades 2 y 3 se combinan de forma clásica: la vulnerabilidad 2 (falta de
sanitización) es la que \textbf{habilita} el XSS, y la vulnerabilidad 3 (cookies sin
protección adicional contra acceso por JavaScript) es lo que hace que ese XSS pueda
\textbf{robar la sesión} concretamente (tal como dice el material sobre XSS: ``permite
robar la sesión''). La vulnerabilidad 1 es un patrón análogo al XXE del material
(el servidor hace un request a donde el atacante decide). La vulnerabilidad 4 conecta con
los principios de diseño generales (no exponer secretos por canales inadecuados), aunque
--- como se avisó arriba --- el material no desarrolla el detalle de headers/OAuth/JWT.

\paragraph{Resolución.}

\textbf{a) Puntos 2 y 3 --- cómo un atacante se suma dinero de forma maliciosa.}
Como el contenido de una noticia se postea \textbf{sin ninguna validación} (vulnerabilidad
2), un publicador (o cualquiera con acceso a publicar) puede subir una noticia cuyo
``contenido'' incluya un \texttt{<script>} malicioso: es un caso de \textbf{XSS
almacenado} (el material define XSS como inyección de JS que se ejecuta en el browser de
la víctima al renderizar HTML sin escapar un dato controlado por el atacante). Cuando
otro usuario visita esa noticia, el JS inyectado se ejecuta en su browser, con la
confianza de la sesión de ese usuario en el sitio.

Como además las cookies de sesión \textbf{no son \texttt{http\_only}} (vulnerabilidad 3),
ese JavaScript malicioso \textbf{puede leer la cookie de sesión} de cualquier usuario que
visite la noticia (incluida la de otro publicador o incluso de un administrador) y
enviarla al atacante --- exactamente la consecuencia que el material atribuye a XSS:
``permite robar la sesión''. Con la cookie de sesión robada, el atacante puede
\textbf{hacerse pasar por ese usuario} (\emph{masquerading}) y usar su identidad
autenticada, por ejemplo, para hacer que el sistema le atribuya a sí mismo (o a una
cuenta propia) visitas/ratings sobre sus propias noticias, o directamente actuar en el
backend con los privilegios de publicador robado. Dado que el modelo de negocio paga
según \textbf{cantidad de visitas}, el atacante puede automatizar, con las sesiones
robadas, visitas/ratings positivos masivos y falsos sobre sus propias noticias
publicadas, generando así ingresos por publicidad de forma fraudulenta sin que el pago
esté respaldado por tráfico/interés real.

\textbf{b) Ventaja que un atacante saca de la vulnerabilidad 1 (SSRF).}
El endpoint \texttt{/\{news\_id\}/positive} (o \texttt{/\{news\_id\}/positive/negative})
hace que el \textbf{servidor} realice una acción a partir de un parámetro (\texttt{news\_id})
que el atacante controla. Si ese parámetro puede manipularse para que en lugar de
identificar una noticia el servidor termine haciendo un request HTTP hacia una URL/host
arbitrario elegido por el atacante, se está en el mismo patrón que el material describe
para \textbf{XXE}: el servidor actúa como intermediario y hace peticiones a destinos que
el atacante no podría alcanzar directamente. Un atacante puede aprovechar esto para:
hacer que el servidor (que sí tiene acceso) haga requests hacia \textbf{recursos internos
de la red} que no son alcanzables desde afuera (por ejemplo, otros servicios internos, o
un endpoint de otra noticia/recurso que solo el backend puede alcanzar), usando al
servidor de noticias como ``proxy'' para escanear o alcanzar esa red interna --- igual
que en el ejemplo de XXE del material, donde el atacante ``puede escanear la red interna
visible por el servidor''. También puede usarlo para forzar al servidor a hacer pedidos
repetidos con fines de denegación de servicio hacia un tercero, usando la infraestructura
del sitio como origen del ataque.

\textbf{c) Por qué es inseguro poner la API KEY en la URL (\texttt{/api/\{API\_KEY\}/publish}) y cómo corregirlo.}
Es inseguro porque la API KEY, al viajar como parte de la \textbf{URL}, deja de ser un
secreto protegido y queda expuesta por canales que no están pensados para llevar
secretos: cualquiera que pueda ver esa URL (quien mire por encima del hombro, quien tenga
acceso a un proxy intermedio, o quien pueda ver dónde queda registrada) accede
directamente a la API KEY con la que se pueden publicar noticias en nombre de ese
publicador. \textbf{Nota:} el material no detalla explícitamente todos los canales por
los que una URL con secreto puede filtrarse (logs, historial del navegador, header
\texttt{Referer}, cachés intermedias); esto excede lo cubierto en las clases provistas,
aunque es consistente con el principio general de que un mecanismo de autenticación
\textbf{no debería fallar de forma insegura} exponiendo la credencial por un canal que no
la protege.

Además, tal como está planteado, no hay forma de saber si quien usa la API KEY es
realmente el publicador dueño de esa clave: cualquiera que la obtenga puede usarla sin
que el sistema pueda distinguir el uso legítimo del malicioso.

\textbf{Corrección propuesta:} no incluir el secreto en la URL; enviar la API KEY en un
\textbf{header HTTP} dedicado a autenticación (p.\ ej. \texttt{Authorization}), que no
queda expuesto por los mismos canales que la URL. Esto es consistente con la idea general
de la materia de separar el canal de datos de control (autenticación) del canal de datos
de la petición en sí, en vez de mezclarlos en el mismo campo que identifica el recurso.

\clearpage
%=====================================================================
\section{Final --- 2025, Segundo Cuatrimestre (1)}
%=====================================================================

%=====================================================================
\subsection{Ejercicio 1 --- Pulseras Bluetooth (AES-CTR, modo degradado)}
%=====================================================================

\begin{consigna}
Eran pulseras bluetooth para abrir puertas que cada dos segundos mandaban su
$(ID, CTR, C)$ donde $CTR$ es un contador, $C = \mathrm{AES\text{-}CTR}_k(CTR \,\|\,
\mathit{bit\_wearing})$. $\mathit{bit\_wearing}$ es un bit que indica si es o no una
persona (si la pulsera está puesta).

El servidor, al recibir $(ID, CTR, C)$, descifra $C$ con la $K$ asignada a ese $ID$
y compara si el $CTR$ descifrado es el mismo que el que venía en claro; si lo es, le
da el acceso.

Si el servidor se queda sin conexión, se entra en un \textbf{estado degradado} donde
la puerta deja pasar a cualquier $ID$ que haya tenido acceso recientemente, sin
hacer ninguna otra validación.

\begin{enumerate}
  \item[a)] ¿Qué validaciones podría hacer el servidor para evitar comportamientos
  extraños como que se abran dos puertas distantes casi al mismo tiempo, o que se
  intenten abrir puertas fuera de un horario normal?
  \item[b)] Si alguien captura los mensajes, explicar cómo podría utilizarlos o
  modificarlos para tener acceso. ¿Cómo lo solucionaría?
  \item[c)] Si consiguen el estado degradado, ¿qué podría hacerse para que no puedan
  acceder repitiendo mensajes?
\end{enumerate}
\end{consigna}

\paragraph{Temas.}
Primer Parte, Clase 2 --- Cifrado (\S Modo CTR: keystream $E_k(\text{nonce}\|i)$,
nunca usa $D_k$, paralelizable, CPA-Secure solo si no se repite $(k,\text{nonce})$);
Primer Parte, Clase 3 --- MACs y Cifrado Autenticado (\S Maleabilidad de cifrados de
flujo/CTR; \S CCA y por qué el cifrado solo no alcanza para integridad; \S MAC:
infalsificabilidad, procesar todo el mensaje); Segunda Parte, Clase 2 ---
Autenticación (\S OTP/HOTP: contador sincronizado, códigos de un solo uso,
throttling); Segunda Parte, Clase 5 --- Seguridad de la Empresa (\S IDS/IPS:
detección por anomalías; \S SIEM: correlación de eventos y alertas).

\paragraph{Qué necesitás saber.}
\begin{itemize}
  \item \textbf{CTR es un cifrado de flujo, no autenticado.} $C_i = E_k(\text{nonce}
  \|i)\oplus M_i$; el descifrado es $M_i = E_k(\text{nonce}\|i)\oplus C_i$ (nunca se
  usa $D_k$). CTR da \emph{confidencialidad} (si no se repite el par clave/nonce),
  pero \textbf{no} da integridad ni autenticación por sí solo.
  \item \textbf{Maleabilidad de cifrados de flujo (Clase 3).} Si
  $\mathrm{Enc}_k(m) = G(k)\oplus m = c$, entonces $c' = c\oplus x \Rightarrow
  \mathrm{Dec}_k(c') = m\oplus x$: un atacante que conoce (o adivina) un bit de
  plano en una posición conocida puede flippear ese bit del cifrado, y al descifrar
  cambia exactamente ese bit del plano, \textbf{sin conocer la clave}. Es el mismo
  mecanismo que el ataque al sueldo en RR.HH.\ del material.
  \item \textbf{Ningún criptosistema de los vistos es CCA-Secure}: el cifrado está
  pensado para confidencialidad, no para resistir manipulación del texto cifrado.
  Por eso la integridad requiere una primitiva aparte: el \textbf{MAC}. Un MAC
  robusto siempre procesa la \emph{totalidad} del mensaje (si no, lo que queda
  afuera se puede alterar sin invalidar la etiqueta).
  \item \textbf{OTP/HOTP (Clase 2, Segunda Parte):} $A$ y $B$ comparten una clave
  $K$ y un contador $C$; el código depende de $(K,C)$ y se recomienda
  \emph{throttling} y (para HOTP) una ventana de \emph{look-ahead} para tolerar
  desincronización del contador.
  \item \textbf{IDS/IPS y SIEM (Clase 5, Segunda Parte):} un IDS/IPS detecta por
  \textbf{anomalías} (además de firmas) y decide bloquear o alertar; un SIEM
  \textbf{correlaciona eventos} de distintas fuentes (aquí, aperturas de puertas
  distintas) para detectar patrones anómalos y generar alertas.
\end{itemize}

\paragraph{Notas y conexiones.}
El enunciado describe un esquema de \emph{confidencialidad sin integridad}: se cifra
$(CTR\|\mathit{bit\_wearing})$ con AES-CTR, pero el mensaje completo $(ID,CTR,C)$ no
lleva ningún MAC que lo ate. Esto es exactamente el escenario que motiva la Clase 3
(``el cifrado por sí solo no protege contra modificaciones''): acá, además, el
$CTR$ viaja \textbf{en claro} junto al cifrado, lo que facilita un replay directo. La
verificación descripta en el enunciado (``compara si el CTR es el mismo que el que
venía en $C$'') es débil porque no dice qué pasa si el $CTR$ ya fue usado antes
(no hay noción de \emph{frescura/monotonicidad}), y no ata $ID$ al contenido cifrado
con una clave de integridad.

\paragraph{Resolución.}

\textbf{a) Validaciones de contexto (correlación entre eventos, no solo por mensaje).}
Estas validaciones no dependen de la criptografía del mensaje individual sino de
\textbf{cruzar eventos en el servidor} (tipo SIEM: correlación de eventos + reglas
de anomalía, IDS/IPS basado en anomalías):
\begin{itemize}
  \item \textbf{Verificación de coherencia física (velocidad de desplazamiento):} el
  servidor conoce la posición de cada puerta. Si el mismo $ID$ abrió una puerta y
  poco después (menos tiempo del físicamente necesario para recorrer la distancia)
  intenta abrir otra puerta lejana, la solicitud es \textbf{físicamente imposible}
  y debe rechazarse y generar una alerta (correlación de eventos por $ID$ a través
  del tiempo y el espacio).
  \item \textbf{Ventanas horarias (ACL temporal):} asociar a cada $ID$ (o a cada
  puerta) un rango horario habilitado en una base de datos; si llega una trama fuera
  de ese rango, denegar el acceso y generar una alerta de seguridad, en lugar de
  autorizar silenciosamente.
  \item En ambos casos el mecanismo es \textbf{detectivo/correlativo} (como un
  IDS/SIEM): no reemplaza la validación criptográfica del mensaje individual, la
  complementa detectando patrones anómalos entre mensajes.
\end{itemize}

\textbf{b) Captura de mensajes: uso/modificación y solución.}
Como $C$ se obtiene con AES-CTR (un cifrado de flujo), aplica directamente la
\textbf{maleabilidad de los cifrados de flujo} vista en la Clase 3:
\begin{itemize}
  \item \textbf{Replay puro:} el atacante captura un mensaje válido $(ID, CTR, C)$ y
  lo reenvía más tarde. El enunciado solo dice que el servidor ``compara si el CTR
  es el mismo que el que viene en $C$'' --- \emph{no} dice que el servidor lleve un
  registro del último $CTR$ usado por ese $ID$. Si no lo lleva, el mensaje capturado
  \textbf{se vuelve a aceptar tal cual} (replay exitoso).
  \item \textbf{Modificación por maleabilidad (si el atacante conoce/adivina el bit
  de \emph{wearing}):} dado que CTR es un XOR con el keystream
  ($C = E_k(\text{nonce}\|i) \oplus (CTR\|\mathit{bit\_wearing})$), sin conocer la
  clave $k$ el atacante puede flippear el bit del cifrado que corresponde a la
  posición del bit $\mathit{bit\_wearing}$ dentro de $C$: eso cambia exactamente ese
  bit al descifrar (igual que el ataque al sueldo: $c'' = c\oplus x$ cambia el bit
  correspondiente del plano), \textbf{sin necesitar la clave}. Así podría forzar
  a que el servidor interprete que la pulsera está puesta (o no) según convenga
  para pasar alguna validación adicional que dependa de ese bit.
  \item \textbf{Solución:} agregar \textbf{integridad} con un MAC sobre el mensaje
  completo $(ID\|CTR\|\mathit{bit\_wearing})$, calculado con una clave de MAC
  (idealmente distinta de la de cifrado) y verificado antes de aceptar el mensaje.
  Un MAC que cubre la totalidad del contenido evita tanto la modificación por
  maleabilidad (cualquier bit flippeado invalida la etiqueta) como facilita atar el
  $ID$ al contenido cifrado. Además, el servidor debe llevar un \textbf{registro del
  último $CTR$ aceptado por $ID$} y solo aceptar $CTR_{\text{nuevo}} > CTR_{\text{último}}$
  (monotonicidad / frescura), lo que bloquea el replay puro.
\end{itemize}

\textbf{c) Estado degradado: evitar replay sin conexión.}
En el modo degradado el servidor no puede verificar nada dinámicamente (no tiene
conexión para consultar estado), así que la defensa tiene que apoyarse en algo que
\textbf{se consuma} localmente, análogo a la idea de \textbf{OTP/HOTP} (Clase 2,
Segunda Parte): claves de un solo uso derivadas de un contador compartido.
\begin{itemize}
  \item Precargar en la memoria de la puerta (mientras hay conexión) una
  \textbf{lista de códigos de un solo uso} (equivalente a un HOTP: valores derivados
  de $(K, C)$ con $C$ incremental), vinculados a cada pulsera/$ID$.
  \item Cuando no hay conexión, la pulsera envía el \textbf{siguiente código no usado}
  de esa lista y la puerta lo valida \textbf{localmente} y lo \textbf{marca como
  consumido} (borrándolo o avanzando un puntero), de forma que no pueda reutilizarse:
  esto es exactamente la garantía de un esquema de un solo uso (OTP), que rompe el
  replay porque cada código sirve una única vez.
  \item Al recuperar la conexión, el servidor debe re-sincronizar el estado
  (equivalente al \emph{look-ahead} de HOTP) y reponer/renovar la lista de códigos
  consumidos localmente.
\end{itemize}

%=====================================================================
\subsection{Ejercicio 2 --- Protocolo de autenticación Cliente--Servidor ($H(K_c\|r)$)}
%=====================================================================

\begin{consigna}
Protocolo de autenticación entre un Cliente $C$ y un Servidor $S$ que ambos tienen
una $K_c$ compartida. La función $H$ es un hash sin colisiones.
\begin{align*}
C \to S &: ID\\
S \to C &: r \quad (\text{número aleatorio})\\
C \to S &: H(K_c \| r)
\end{align*}
El servidor calcula el hash usando la $K_c$ que tiene guardada y el $r$. Si da lo
mismo que recibió, autentica a $C$.

Preguntas:
\begin{enumerate}
  \item[1.] ¿Por qué el servidor manda ese $r$? ¿Qué pasaría si el cliente solo
  mandase $H(K_c)$ en vez de que se use el $r$?
  \item[2.] ¿Disminuiría la seguridad si se usara un contador incremental en vez
  del $r$ aleatorio?
  \item[3.] Plantear cómo se podría hacer para que el cliente autentique también al
  servidor, sin introducir nuevas claves.
\end{enumerate}
\end{consigna}

\paragraph{Temas.}
Segunda Parte, Clase 2 --- Autenticación (\S Challenge-Response y EKE: el desafío
$r$ como nonce, ``el secreto nunca viaja por la red''); Primer Parte, Clase 5 ---
Protocolos (\S Ataque por reflexión y ``romper la simetría'' con nonces/identidad
por sentido; \S Handshakes con nonces de ambos lados, ej.\ TLS \emph{Hello} con
$r_1, r_2$).

\paragraph{Qué necesitás saber.}
\begin{itemize}
  \item \textbf{Challenge-response (Clase 2, Segunda Parte):} el servidor manda un
  \textbf{challenge único} $r$ (nonce); el cliente responde $f(k,r)$ (acá,
  $H(K_c\|r)$) probando que conoce el secreto $K_c$ \textbf{sin enviarlo}. El
  diagrama del material remarca explícitamente: ``el nonce evita reutilizar
  respuestas capturadas''.
  \item Si no hubiera $r$ (o el cliente mandara siempre $H(K_c)$, un valor fijo), la
  respuesta sería \textbf{siempre la misma}: un atacante que la capture una vez
  puede \textbf{reenviarla} (replay) y autenticarse sin conocer $K_c$.
  \item Un \textbf{contador incremental} en vez de un $r$ aleatorio también da
  frescura (cada valor se usa una sola vez si el servidor lo verifica creciente),
  pero es \textbf{predecible}: un atacante que observa la secuencia puede anticipar
  el próximo valor. En el material esto es exactamente el punto que distingue HOTP
  (contador, requiere sincronización, predecible en su progresión) de un valor
  \textbf{aleatorio} imprevisible.
  \item \textbf{Autenticación mutua sin nuevas claves:} el material resuelve el
  problema simétrico en los \textbf{handshakes} (Clase 5, Primer Parte) generando
  un desafío \textbf{en cada sentido} con la misma clave compartida, y remarca
  (a propósito del ataque por reflexión) que hay que \textbf{romper la simetría}
  entre ambos desafíos (por ejemplo, atando la identidad/dirección a cada
  desafío) para que una sesión paralela no pueda usarse para reflejar la
  respuesta.
\end{itemize}

\paragraph{Notas y conexiones.}
Este protocolo es el patrón \textbf{challenge-response} de la Clase 2 (Segunda
Parte) aplicado con hash en lugar de una función $f(k,r)$ genérica. El ataque por
\textbf{reflexión} de la Clase 5 (Primer Parte) es la advertencia clave para la
pregunta 3: si simplemente se agrega un segundo desafío simétrico (el cliente
manda un $r_c$ y espera $H(K_c\|r_c)$ del servidor, de la misma forma en que el
servidor se lo pide al cliente), un atacante podría abrir una \textbf{sesión
paralela} y reflejar el desafío recibido en una sesión para responder en la otra,
si no se distingue el rol/dirección de cada desafío.

\paragraph{Resolución.}

\textbf{1) Por qué se manda $r$; qué pasa con $H(K_c)$ solo.}
El servidor manda $r$ para dar \textbf{frescura} a la respuesta: como $r$ es
distinto en cada intento de autenticación, la respuesta $H(K_c\|r)$ también lo es,
así que \textbf{capturar una respuesta pasada no sirve para autenticarse de
nuevo} (no hay dos ejecuciones del protocolo con la misma respuesta válida). Si
el cliente mandara directamente $H(K_c)$ (sin usar $r$), la respuesta sería
\textbf{siempre el mismo valor fijo} en cada autenticación; un atacante que la
capture una sola vez (escuchando el canal) puede luego \textbf{reenviarla
(replay)} y el servidor la aceptaría como si fuera el cliente legítimo, sin
conocer nunca $K_c$. Es exactamente el rol del nonce en challenge-response: ``el
secreto nunca viaja por la red'' y \textbf{cada respuesta es única}, ligada a un
desafío que no se repite.

\textbf{2) Contador incremental en vez de $r$ aleatorio.}
Sí \textbf{disminuye la seguridad} respecto de un nonce verdaderamente aleatorio,
aunque igual aportaría algo de frescura si el servidor exige que sea creciente
(evitaría el replay literal del mismo mensaje). El problema es la
\textbf{predictibilidad}: si el contador es incremental y el atacante observa la
secuencia (por ejemplo, escuchando varias rondas), puede \textbf{anticipar el
valor del próximo contador} y, si además logra capturar u obtener de otra forma el
valor de $H(K_c\|r)$ para un $r$ futuro conocido (o construir un ataque que
aproveche esa previsibilidad, según qué tan expuesto esté el mecanismo), reduce el
espacio de incertidumbre del protocolo. En el material esto es análogo a la
comparación entre HOTP (contador, requiere sincronización explícita y su
progresión es previsible) y un desafío aleatorio (TOTP/nonce, imprevisible salvo
que se conozca la clave). Un valor \textbf{aleatorio e imprevisible} es preferible
porque no permite a un atacante preparar de antemano una respuesta para un
desafío futuro.

\textbf{3) Autenticación mutua sin introducir nuevas claves.}
Usando la misma $K_c$ compartida, se agrega un desafío \textbf{en el otro
sentido}, rompiendo la simetría para evitar un ataque por reflexión (Clase 5,
Primer Parte):
\begin{align*}
C \to S &: r_C,\ ID\\
S \to C &: r_S,\ H(K_c \| r_C)\\
C \to S &: H(K_c \| r_S)
\end{align*}
El servidor prueba que conoce $K_c$ respondiendo al desafío $r_C$ del cliente
(en el mismo mensaje en que manda su propio desafío $r_S$); el cliente, a su vez,
responde al desafío $r_S$. Cada parte verifica el hash correspondiente a \emph{su
propio} desafío. Es importante que ambos desafíos \textbf{no sean intercambiables
uno por el otro} sin ambigüedad (en el material, la defensa contra la reflexión es
distinguir la dirección/rol de cada nonce, por ejemplo incluyendo dentro del hash
la identidad o el rol del que responde, del estilo $H(K_c\|r_S\|\text{``cliente''})$
y $H(K_c\|r_C\|\text{``servidor''})$) para que una sesión paralela abierta por un
atacante no pueda usarse para que el servidor legítimo ``firme'' un desafío que
luego se reutilice contra el cliente. No se introduce ninguna clave nueva: se
reutiliza $K_c$ y $H$ en ambos sentidos.

%=====================================================================
\subsection{Ejercicio 3 --- CasaIA (IoT hogareño): hipótesis de falla}
%=====================================================================

\begin{consigna}
Una página web y aplicación mobile conectadas con una API REST. Se llamaba CasaIA,
te dejaba controlar los dispositivos de tu casa a través de wifi, y podrías
configurar funciones por horario y por ubicación geográfica del teléfono. Los
dispositivos están conectados con un websocket en TCP al servidor, que es el que
les puede mandar comandos (abrir/cerrar, etc.).

Los dispositivos se vinculan escaneando un QR que viene en ellos.

Se puede dar acceso temporal a través de la página a otra persona: el usuario pone
el mail de la persona y a esta le llega un mail con el cual, al abrirlo, consigue
acceso temporal. También se genera un token de sesión al lograrse el acceso.

Se pide: como siempre, dar \textbf{4 hipótesis de falla} y su prueba.
\end{consigna}

\paragraph{Temas.}
Segunda Parte, Clase 6 --- Pentesting (\S Metodología de hipótesis de falla, Paso 2:
áreas candidatas de hipótesis; \S la prueba debe ser repetible y concluyente);
Segunda Parte, Clase 3 --- Principios de Diseño y Vulnerabilidades (\S CSRF; \S XSS;
\S nonce/token anti-replay); Segunda Parte, Clase 2 --- Autenticación (\S OTP/HOTP,
contador vs.\ aleatorio; \S challenge-response); glosario (\S OAuth2/tokens,
\S nonce).

\paragraph{Qué necesitás saber.}
\begin{itemize}
  \item \textbf{Metodología de hipótesis de falla (Clase 6):} el paso 2 consiste en
  plantear posibles vulnerabilidades \textbf{examinando implementaciones y
  mecanismos} (¿está bien implementado el control?, ¿el ambiente introduce
  errores?, ¿es realmente seguro?) y \textbf{comparando con sistemas parecidos}
  (sistemas similares tienen problemas similares). El paso 3 (prueba) exige que la
  verificación de la hipótesis sea \textbf{repetible y concluyente}.
  \item \textbf{Nonce / frescura (Clase 3):} la defensa real contra CSRF es un
  \textbf{token/nonce único} que dé frescura y evite replay; lo mismo aplica a
  cualquier enlace o token de un solo uso (como el mail de acceso temporal de
  CasaIA).
  \item \textbf{XSS (Clase 3):} inyectar código que termine ejecutado por el
  browser de la víctima permite \textbf{robar la sesión} (aprovechando la confianza
  del sitio en su sesión/cookie).
  \item \textbf{Tokens de sesión (glosario):} un token de sesión, si no tiene
  expiración adecuada, alcance limitado, o viaja sin protección, puede ser robado o
  reutilizado para suplantar al usuario.
  \item \textbf{OTP/HOTP (Clase 2):} un código de acceso temporal enviado por mail
  que actúa como un ``one-time'' debería, como en HOTP/TOTP, tener \textbf{throttling}
  y no ser reutilizable ni de vida indefinida.
\end{itemize}

\paragraph{Nota sobre alcance (regla de oro).}
El enunciado pide una \textbf{metodología puntual} (``4 hipótesis de falla, una
prueba cada una''); en el material de Clase 6 (Pentesting) la metodología está
descripta en \textbf{5 pasos} (Recolección $\to$ Hipótesis $\to$ Prueba $\to$
Generalización $\to$ Eliminación) y el paso de \textbf{Hipótesis} es explícitamente
donde se listan posibles vulnerabilidades (no hay, en las clases provistas, una
regla de puntaje del tipo ``4 hipótesis, 1 punto cada una''; eso parece ser una
convención de la consigna del parcial, no un contenido de la clase). Lo que sí está
en el material es \emph{cómo} se plantea y se prueba cada hipótesis (paso 2 y 3 de
la metodología), que es lo que se aplica a continuación sobre CasaIA.

\paragraph{Resolución.}
Aplicando la metodología de hipótesis de falla (examinar implementación/mecanismos
+ comparar con sistemas parecidos) a los cuatro componentes distintivos de CasaIA
(QR de vinculación, websocket de comandos, acceso temporal por mail, token de
sesión):

\begin{enumerate}
  \item \textbf{Hipótesis --- vinculación por QR sin verificación de posesión.} Si
  el QR del dispositivo solo contiene un identificador estático (sin un secreto o
  desafío que se consuma una única vez), cualquiera que \textbf{fotografíe o copie}
  el QR (por ejemplo, publicado sin querer en una foto, o visto por un instalador)
  podría vincular ese dispositivo a su propia cuenta. \emph{Prueba (repetible y
  concluyente, Clase 6):} tomar una foto del QR de un dispositivo ya vinculado y,
  desde otra cuenta, intentar vincularlo nuevamente; si el sistema lo permite (o no
  invalida la vinculación anterior), la hipótesis queda confirmada.
  \item \textbf{Hipótesis --- falta de autenticación/autorización en los mensajes
  del websocket.} Si el canal websocket entre servidor y dispositivo no reautentica
  cada comando (solo se autenticó la conexión TCP inicial) ni valida que el comando
  provenga de un usuario autorizado sobre \emph{ese} dispositivo puntual, un
  atacante que logre interceptar o inyectar tráfico en ese canal podría mandar
  comandos (abrir/cerrar) sin ser el dueño. \emph{Prueba:} con acceso a la red del
  dispositivo (caja gris), capturar el tráfico del websocket y reenviar
  (\emph{replay}) un comando ``abrir'' capturado, o inyectar un comando propio, y
  verificar si el servidor lo ejecuta sin re-validar autorización.
  \item \textbf{Hipótesis --- el enlace de acceso temporal por mail no tiene
  frescura/expiración adecuada (análogo a CSRF/replay, Clase 3).} Si el link que
  llega por mail para dar acceso temporal es un token que no expira, no es de un
  solo uso, o es predecible/corto, cualquiera que lo intercepte (compartido reenvío
  de mail, mail comprometido, o el link cacheado en el navegador) podría usarlo
  para obtener acceso indefinidamente, igual que un token anti-CSRF débil no da
  frescura real. \emph{Prueba:} generar un acceso temporal, usarlo una vez, y
  volver a usar el mismo enlace pasado un tiempo prudencial (o tras revocar el
  acceso desde la app); si sigue funcionando, la hipótesis queda confirmada.
  \item \textbf{Hipótesis --- robo/reutilización del token de sesión (estilo
  XSS/robo de sesión, Clase 3).} Si el token de sesión generado al lograrse el
  acceso (temporal o del dueño) no tiene alcance ni expiración acotados, o si la
  interfaz web es vulnerable a inyección de código (XSS) que permita leer ese
  token, un atacante podría robarlo y \textbf{suplantar la sesión} sin conocer
  ninguna contraseña. \emph{Prueba:} intentar inyectar un script en algún campo de
  la aplicación web (por ejemplo, el nombre del dispositivo o de la persona
  invitada) y verificar si se ejecuta al ser renderizado por otro usuario
  (confirmando que el token/cookie de sesión de ese usuario quedaría expuesto);
  alternativamente, verificar si el token de sesión sigue siendo válido después de
  que expiró el acceso temporal que lo originó.
\end{enumerate}


\clearpage
%=====================================================================
\section{Ejercicios adicionales}
%=====================================================================

% =====================================================================
\subsection{Ejercicio --- Bingo virtual}
% =====================================================================

\begin{consigna}
Una franquicia de bingos decide crear una aplicación que simule un bingo virtual donde
los participantes paguen y obtengan dinero real.

Cada juego tiene un valor de ingreso, y le permite al participante obtener un cartón con
15 número elegidos al azar entre 1 y 90. Cuando comienza el juego, el servidor comienza a
``sacar'' números, y los participantes deben ir marcando los números en su cartón si
corresponde.

Cuando uno de los participantes completa los 15 número de su cartón, selecciona el botón
``bingo'', y si realmente salieron todos, gana el premio de esa jugada y termina la
partida para todos. Si no fuese así, se le avisa y se lo descalifica.

Resuelta la seguridad de los aspectos de introducir y retirar dinero se identificaron
tres funcionalidades necesarias respecto a la dinámica del juego. Explicar como se podría
lograr cada una recurriendo a funciones criptográficas \underline{simétricas solamente},
justificando la elección. De ser necesario, explicar que información se intercambia en
cada momento entre las apps y el servidor.

\begin{enumerate}
  \item[a)] El servidor debe generar una secuencia de números entre 1 y 90, que luego
  pueda ser reproducida por los clientes al finalizar la partida (para garantizar que no
  se favoreció adrede a alguien en el sorteo). (1 pto.)
  \item[b)] Los cartones se crean en las apps cliente al inicio de la partida. El servidor
  no conoce los cartones hasta el momento de verificar un bingo, pero debe ser posible
  para el servidor verificar que el cartón no fue modificado una vez empezado el juego
  (para evitar que un jugador arme el cartón ideal para ganar luego de 15 números). (1 pto.)
  \item[c)] La app de cada participante deberá poder verificar si el ganador realmente
  completó el bingo y emitir una alerta en caso contrario (para evitar que el servidor de
  por ganador a alguien que no completo el cartón). Las apps no se comunican entre si,
  sino a través del servidor. (1 pto.)
\end{enumerate}

\emph{Nota}: Considerar que hay aproximadamente $2^{55}$ cartones (combinatoria de 90
tomados de a 15), y asumir que un atacante podría llegar a probar esa cantidad de
combinaciones.
\end{consigna}

\paragraph{Temas.}
Clase 3 --- MACs y Cifrado Autenticado (\S Qué es un MAC; \S HMAC; \S funciones de hash y
sus propiedades de resistencia; \S usos de la resistencia a preimagen para protocolos de
compromiso/prueba de conocimiento; \S no-repudio: por qué los MAC no lo dan).

\paragraph{Qué necesitás saber.}
\begin{itemize}
  \item \textbf{MAC/HMAC:} terna $(\mathrm{Gen},\mathrm{Mac},\mathrm{Vrfy})$ con clave
  \textbf{compartida} $k$; da integridad y autenticación de origen \emph{entre partes que
  comparten la clave}, pero \textbf{no} no-repudio.
  \item \textbf{Función de hash:} sin clave, de tamaño fijo, con tres resistencias
  (preimagen, segunda preimagen, colisión). La resistencia a preimagen es la base de los
  \textbf{protocolos de prueba de conocimiento / compromiso}: publicar $H(x)$ y después
  revelar $x$, de forma que nadie pudo deducir $x$ antes de la revelación, y quien reveló
  $x$ no puede haber cambiado de valor sin que se note (segunda preimagen).
  \item \textbf{Ataque de cumpleaños / brute-force de espacios chicos:} si el espacio de
  posibles mensajes es chico (acá, $2^{55}$ cartones), un compromiso $H(\text{cartón})$
  \emph{sin sal} es vulnerable a que cualquiera enumere los $2^{55}$ cartones posibles,
  los hashee y compare contra el compromiso publicado para deducir el cartón antes de
  tiempo. La nota del enunciado es justamente el gancho para exigir un \textbf{nonce/sal
  grande} (mucho mayor a 55 bits) dentro del compromiso.
\end{itemize}

\paragraph{Notas y conexiones.}
Los tres incisos son variaciones del mismo patrón: (a) y (b) son \textbf{compromisos}
(commitments) usando la resistencia a preimagen/segunda preimagen de una función de hash
(no llevan clave, pero son ``funciones criptográficas simétricas'' en el sentido amplio
que exige la consigna, ya que no dependen de un par de claves pública/privada); (c) es un
esquema de \textbf{autenticación con clave compartida} (MAC/HMAC) para que todas las apps
puedan validar de forma independiente los mismos mensajes del servidor. Como está prohibido
usar asimétrica, no hay firma digital disponible, así que ninguna de las tres soluciones
da no-repudio frente al servidor (el servidor y las apps comparten secretos).

\paragraph{Resolución.}

\textbf{a) Sorteo verificable.} Antes de arrancar la partida el servidor elige una
\textbf{semilla secreta} $s$ (número aleatorio, e.g. 256 bits) y publica un
\textbf{compromiso} $c=H(s)$ a todas las apps (usando la resistencia a preimagen de un
hash: nadie puede deducir $s$ a partir de $c$). La secuencia de 90 números se deriva de
forma \textbf{determinística} de $s$ con una función pseudoaleatoria con clave (p.ej.
$\mathrm{HMAC}_s(i)$ para $i=1,2,\dots$, usada como fuente de aleatoriedad de un
\emph{shuffle} tipo Fisher--Yates sobre $\{1,\dots,90\}$). El servidor va revelando los
números uno a uno a medida que los ``saca''; como se derivan de $s$ que solo él conoce, no
puede ser forzado a revelar otro número, pero tampoco puede cambiar el orden a mitad de
partida sin que se note (porque ya se comprometió a $c=H(s)$ antes de empezar). Al
finalizar la partida, el servidor revela $s$; cada app recomputa $H(s)$ y verifica que
coincide con $c$ (garantiza que $s$ es el mismo que se comprometió al inicio, por
resistencia a segunda preimagen), y recomputa la secuencia completa de 90 números con $s$
para verificar que coincide con la secuencia efectivamente anunciada durante el juego. Si
alguna etapa no coincide, la secuencia fue manipulada.

\textbf{b) Cartón no modificable.} Al empezar la partida, cada app cliente genera su
cartón y envía al servidor un \textbf{compromiso} $h=H(\text{cartón}\,\|\,r)$, donde $r$ es
un \textbf{nonce aleatorio grande} (mucho mayor a 55 bits, p.ej. 128 bits) generado por la
app y que \textbf{no} se envía al servidor en ese momento. Por la nota del enunciado, si se
omitiera $r$ el compromiso $H(\text{cartón})$ sería vulnerable a que cualquiera (el propio
servidor u otro atacante con el compromiso) recorra los $2^{55}$ cartones posibles,
hasheándolos hasta encontrar el que coincide, revelando el cartón antes de que termine la
partida. Agregar $r$ (con un espacio mucho mayor que $2^{55}$) hace que la búsqueda por
fuerza bruta sea inviable ($2^{55+128}$ combinaciones), preservando la \emph{ocultación}
del cartón hasta el final. Al momento de reclamar bingo (o al finalizar la partida), la app
envía $(\text{cartón}, r)$ al servidor, que recalcula $H(\text{cartón}\,\|\,r)$ y verifica
que coincide con el $h$ recibido al inicio (resistencia a segunda preimagen: no puede
existir otro cartón$'$ con el mismo $h$, así que el jugador no pudo cambiar el cartón
después de comprometerse).

\textbf{c) Verificación distribuida sin comunicación entre apps.} Todas las apps
comparten con el servidor (instalada de antemano, al registrarse) una \textbf{clave
simétrica} $k$ (puede ser una por partida, distribuida al inicio junto con $(G,q,g)$... o
más simple, generada por el servidor y enviada cifrada/instalada a cada app al arrancar la
partida). Cada número que el servidor ``saca'' en el paso (a) se transmite junto con una
etiqueta $t_i=\mathrm{HMAC}_k(i\,\|\,\text{número}_i)$; cada app, al recibirlo, valida
$\mathrm{Vrfy}_k$ y solo si es válido lo marca como ``anunciado''. De esta forma todas las
apps arman localmente el mismo historial \textbf{autenticado} de números sacados, sin
hablar entre sí (solo escuchan al servidor). Cuando alguien grita bingo, el servidor
retransmite a todas las apps el cartón revelado del ganador junto con $(r,h)$ del punto
(b); cada app, de forma independiente, (i) recalcula $H(\text{cartón}\|r)$ y compara con el
$h$ que el servidor había repartido al inicio de la partida (si no coincide, el cartón fue
alterado), y (ii) chequea que los 15 números del cartón estén todos entre los que su propio
historial local marcó como ``anunciados y autenticados con $k$''. Si cualquiera de las dos
verificaciones falla, la app emite una alerta al servidor (sigue sin comunicarse
directamente con las demás apps, tal como pide la consigna).

\begin{ojo}
Ninguna de las tres soluciones da \textbf{no-repudio}: como todo se basa en hash sin clave
o en un secreto \textbf{compartido} ($s$, $r$, $k$), el servidor sigue siendo, en última
instancia, quien controla la información y a quien hay que confiarle la ejecución honesta
del protocolo (el enunciado no permite usar firma digital, que sería la única forma de dar
no-repudio, según Clase 4).
\end{ojo}

% =====================================================================
\subsection{Ejercicio --- Robo de cables eléctricos}
% =====================================================================

\begin{consigna}
El robo de cables por el valor intrínseco del cobre que traen es un gran problema para las
compañías de distribución eléctrica, dado que les genera un doble perjuicio: el costo
económico de reponer los cables y el costo asociado a resarcir a clientes cuyo servicio se
ve afectado.

Para detectar cortes (intencionales o no) en las lineas, cada cierta distancia se conectan
pequeños dispositivos que modulan una señal de alta frecuencia en la misma linea que
transporta energía. Cada dispositivo tiene una firma única (una frecuencia y un periodo
dado --- diseñado de tal forma que todos los dispositivos pueden inyectar su señal sin
interferirse entre si). Las lineas terminan en transformadores. Antes del transformador se
coloca un dispositivo llamado \emph{colector} que recolecta todas las señales y aplica un
filtro de alta frecuencia que limpia la señal para que llegue al transformador sin este
ruido.

Por la forma en que se juntan las señales, cada dispositivo es capaz de transmitir el
equivalente a una secuencia de 320 bits arbitraria.

Una implementación en particular de este sistema asigna a cada dispositivo un par de claves
publica y privada de 320 bits, y hace que cada dispositivo envíe regularmente la firma
digital de su numero de serie único (solo la firma!).

El \emph{colector} conoce los números de serie de todos los dispositivos, y de cada uno la
frecuencia (lo que le permite extraer los 320 bits de la señal especifica), la ubicación
relativa en el cable y su clave publica. De esta manera, puede determinar que si deja de
recibir la señal del dispositivo $x$ y todos los anteriores hubo un corte en el cable,
mientras que si solo falta $x$ hubo un problema con el dispositivo en particular.

\begin{enumerate}
  \item[a)] Explicar paso a paso que tiene que hacer el colector para validar el estado de
  la linea y llegar a uno de tres estados: todo ok, corte cercano al dispositivo $x$,
  problemas con el dispositivo $x$. Ser preciso en la aplicación de cada función
  criptográfica involucrada y sus parámetros. (1 pto)
  \item[b)] En busca de abaratar costos, se probó una modificación al sistema donde
  \underline{todos} los dispositivo comparten una \underline{única} clave secreta con el
  aparato terminal (aunque cada uno sigue teniendo un numero de serie diferente), y
  transmiten una etiqueta de su numero de serie, generada con un MAC criptográficamente
  seguro. ¿Como cambia el algoritmo del colector? ¿Que desventajas genera este cambio a
  nivel seguridad? (1 pto)
  \item[b)] ¿Como modificaría el protocolo para evitar un ataque donde alguien corte una
  sección pero agregue un aparato que repita las señales enviadas previamente por los
  dispositivos que quedaron fuera de linea? (1 pto)
\end{enumerate}

\emph{Nota}: es imposible modificar el ancho de banda de cada dispositivo de 320 bits.
\end{consigna}

\paragraph{Temas.}
Clase 4 --- Cifrado Asimétrico y Firma Digital (\S Firmas digitales; \S RSA-Signature y por
qué se firma el hash; \S curvas elípticas: 320 bits $\approx$ seguridad de RSA-2048);
Clase 3 --- MACs y Cifrado Autenticado (\S Qué es un MAC; \S no-repudio: por qué los MAC no
lo dan); Clase 5 --- Protocolos (\S nonce y timestamp como mecanismos de frescura;
\S ataques de replay y sus defensas).

\paragraph{Qué necesitás saber.}
\begin{itemize}
  \item \textbf{Firma digital:} $\mathrm{Sign}_{sk}(m)$ con la clave \textbf{privada} del
  firmante; $\mathrm{Vrfy}_{pk}(m,\sigma)$ con la clave \textbf{pública}, que cualquiera
  (incluido el colector) puede ejecutar sin compartir secretos. Da no-repudio e integridad,
  y es \textbf{públicamente verificable}: el colector solo necesita la clave pública de
  cada dispositivo, no un secreto compartido.
  \item \textbf{320 bits como firma sobre curva elíptica:} el material de la clase indica
  que para una seguridad equivalente a RSA-2048 ($\sim$128 bits) alcanza con
  $\sim$320 bits en curvas elípticas (ECDSA/ElGamal sobre curva). Esto encaja exactamente
  con el ancho de banda de 320 bits que puede transmitir cada dispositivo: es coherente con
  que la ``firma digital de su número de serie'' quepa en una única transmisión.
  \item \textbf{MAC con clave compartida entre todos:} da integridad/autenticación pero
  \textbf{no} distingue quién generó la etiqueta (cualquiera con la clave puede forjar la
  de cualquier otro) y \textbf{no} da no-repudio (Clase 3).
  \item \textbf{Frescura (nonce/timestamp) contra replay:} Clase 5 muestra que sin un
  mecanismo de frescura, un atacante puede grabar y reenviar (\emph{replay}) un mensaje
  válido viejo. La defensa es incluir un \textbf{nonce} o \textbf{timestamp} en lo que se
  firma/etiqueta, de forma que el verificador rechace mensajes ``viejos''.
\end{itemize}

\paragraph{Notas y conexiones.}
El enunciado arma un caso de uso real de firma digital asimétrica (Clase 4) con la
restricción de ancho de banda de 320 bits, que \textbf{calza justo} con el tamaño de firma
de un esquema sobre curva elíptica visto en la Clase 4. La parte (b) reproduce el
contraste MAC vs.\ firma digital de la Clase 3 (no-repudio) y la parte (b) repetida conecta
con los mecanismos de frescura/anti-replay de la Clase 5 (Needham--Schroeder,
Denning--Sacco), pero con la restricción extra de que el tamaño transmitido no puede
crecer: se resuelve firmando/etiquetando el \textbf{hash} de (serie $\|$ contador), tal
como en Clase 4 se explica por qué se firma el hash y no el mensaje (permite mensajes de
tamaño arbitrario con una firma de tamaño fijo).

\paragraph{Resolución.}

\textbf{a) Algoritmo del colector (esquema original, con firma digital).} El colector
mantiene, para cada dispositivo $i=1,\dots,n$ ordenado por \textbf{ubicación relativa en el
cable} (desde el más cercano al colector hasta el más lejano), su número de serie
$\text{serie}_i$, su frecuencia/periodo (para poder sintonizar y extraer los 320 bits que
le corresponden) y su clave pública $pk_i$. Paso a paso:
\begin{enumerate}
  \item Para cada dispositivo $i$, el colector \textbf{sintoniza} su frecuencia/periodo
  específica y trata de capturar los 320 bits transmitidos en esa señal.
  \item Si capta señal en la frecuencia de $i$: interpreta esos 320 bits como una firma
  $\sigma_i$ y ejecuta $\mathrm{Vrfy}_{pk_i}(\text{serie}_i,\sigma_i)$. Si la verificación
  da $1$, el dispositivo $i$ está \textbf{transmitiendo correctamente} (nadie puede haber
  falsificado esa firma sin la clave privada de $i$).
  \item Si \textbf{no} capta señal en la frecuencia de un dispositivo $x$ y tampoco en la
  de \textbf{todos} los dispositivos ubicados a partir de $x$ hacia el otro extremo del
  cable (los que están ``detrás'' de $x$ respecto del colector): como la señal viaja por
  el mismo conductor de cobre, un corte físico en el punto donde está $x$ interrumpe la
  continuidad eléctrica para todos los dispositivos más allá del corte. Estado:
  \textbf{corte cercano al dispositivo $x$}.
  \item Si \textbf{solo falta} la señal de $x$, pero los dispositivos ubicados antes
  \textbf{y} después de $x$ en el cable sí llegan (y verifican $\mathrm{Vrfy}$ = 1): el
  cable sigue eléctricamente continuo (si no, tampoco llegarían los de más allá de $x$),
  así que la falla es puntual del propio dispositivo $x$ (rotura, batería, etc.). Estado:
  \textbf{problema con el dispositivo $x$}.
  \item Si todos los dispositivos verifican correctamente: estado \textbf{todo ok}.
\end{enumerate}

\textbf{b) Variante con clave secreta única (MAC compartido).} El algoritmo del colector
cambia únicamente en el paso de verificación: en vez de $\mathrm{Vrfy}_{pk_i}$ con una
clave pública distinta por dispositivo, el colector (que ahora debe conocer la
\textbf{única} clave secreta $K$ compartida por todos los dispositivos) recibe la etiqueta
$t_i=\mathrm{MAC}_K(\text{serie}_i)$ transmitida en los 320 bits, recalcula
$\mathrm{MAC}_K(\text{serie}_i)$ y compara con $t_i$ (en vez de verificar una firma con
clave pública, recalcula y compara, como cualquier MAC). El resto de la lógica (corte
cercano a $x$ vs.\ problema puntual en $x$, según qué señales llegan) no cambia.

\textbf{Desventajas de seguridad:}
\begin{itemize}
  \item \textbf{Punto único de falla:} al ser una clave \textbf{compartida por todos los
  dispositivos}, si un atacante extrae físicamente $K$ de \emph{un solo} dispositivo (por
  ejemplo abriéndolo), \textbf{compromete a todo el sistema}: puede generar etiquetas
  válidas para cualquier número de serie, falsificando la señal de cualquier dispositivo
  (incluso de los que nunca tocó).
  \item \textbf{Sin no-repudio ni atribución:} como todos comparten $K$, la etiqueta
  \textbf{no prueba que la transmitió ese dispositivo en particular}: en principio
  cualquiera que conozca $K$ (incluido otro dispositivo legítimo, o el propio colector)
  podría haber generado esa etiqueta. Con firma digital (parte a) cada dispositivo es el
  único que puede producir su firma (clave privada propia); con MAC compartido esa
  propiedad se pierde por completo.
  \item El colector mismo pasa a ser un blanco crítico: si se compromete, se filtra $K$ y
  cae la seguridad de \emph{todos} los dispositivos de la red, no solo de uno.
\end{itemize}

\textbf{b) (repetida) --- Evitar el replay de señales viejas.} El ataque descripto es un
\textbf{replay}: alguien corta una sección y coloca un aparato que retransmite las
firmas/etiquetas \textbf{viejas} capturadas de los dispositivos que quedaron desconectados,
haciendo que el colector siga viendo ``todo ok'' aunque esos dispositivos ya no estén. La
defensa (Clase 5) es agregar \textbf{frescura} (nonce o timestamp/contador) a lo que se
firma o etiqueta. Como el ancho de banda está fijado en 320 bits y no se puede agrandar, no
alcanza con agregar el contador ``en claro'' junto a la firma (ocuparía más bits); la
solución es la misma idea que justifica firmar el \textbf{hash} en la Clase 4: en vez de
firmar/etiquetar $\text{serie}_i$ solo, cada dispositivo firma/etiqueta
$H(\text{serie}_i\,\|\,c_i)$, donde $c_i$ es un \textbf{contador} (o ventana de tiempo)
que el dispositivo incrementa en cada transmisión y que el colector puede predecir
(ambos, dispositivo y colector, mantienen el contador sincronizado, o el colector acepta
una ventana de contadores cercana al último visto). Como el hash tiene salida de tamaño
fijo, la firma/etiqueta sigue ocupando exactamente 320 bits, sin importar que ahora el
mensaje firmado sea más largo. El colector, al verificar, recalcula
$H(\text{serie}_i\|c_i^{\text{esperado}})$ y ejecuta $\mathrm{Vrfy}$ (o recomputa el MAC)
contra ese valor; si el aparato del atacante repite una firma/etiqueta vieja, el contador
que usó para generarla ya \textbf{no coincide} con el esperado por el colector, y la
verificación falla (detecta el replay). Si el contador nunca avanza mientras se supone que
el dispositivo sigue transmitiendo, el colector puede además usar eso como señal adicional
de anomalía.

% =====================================================================
\subsection{Ejercicio --- Organizame (pentesting de una app de calendario)}
% =====================================================================

\begin{consigna}
``Organizame'' es una aplicación móvil que busca potenciar el uso eficiente del tiempo de
las personas.

La aplicación permite que un usuario se registre utilizando su cuenta de Google o su Apple
Id, y pide acceso al calendario de dicha cuenta para obtener información de las reuniones
agendadas.

La información de agenda es almacenada en una base de datos en el servidor, donde procesos
corren regularmente análisis y proponen mejoras a la agenda. Cuando uno de estos procesos
encuentra una posible mejora, envía una notificación a la aplicación para que el usuario
decida si desea aplicar el cambio o no. Algunos procesos son simples y rápidos, aunque hay
algunos que tienen un orden de complejidad $O(n^3)$ en la cantidad de eventos del mes.

La aplicación permite compartir huecos de tiempo libre del calendario de forma publica,
para que otras personas puedan agendar reuniones. Para ello disponibiliza una pagina web
por usuario, de la forma: \texttt{https://organizame.io/\{user\_id\}/public/calendar}. Un
usuario puede habilitar o deshabilitar dicha página en su configuración de perfil. En ella,
cualquier persona puede agendar una reunión en los momentos disponibles, y dejarle al dueño
un mensaje que verá en su calendario.

Ante un pedido de realizar un pentest que abarque el sistema, establecer 4 hipótesis de
falla \textbf{pertinentes} para el caso, basadas en problemas de seguridad en aplicaciones.
Explicar \textbf{para cada una}, como probarla según la metodología (la acción concreta que
ejecutaría para verificar o refutar la hipótesis). (4 puntos, 1 pto cada uno)

\textbf{Aclaración}: No utilizar problemas genéricos como hipótesis de falla. Plantear
situaciones y lugares concretos donde podría estar el problema. Por ejemplo, en lugar de
decir ``puede haber un buffer overflow'', decir ``es probable que haya un buffer overflow
en la lectura del parámetro x que debe interpretarse como número y puede tener un buffer
pequeño asumiendo que el número tiene menos de los dígitos de un long''.
\end{consigna}

\paragraph{Temas.}
Clase 6 (Segunda Parte) --- Pentesting (\S Metodología de hipótesis de falla, en
particular Paso 2 ``Hipótesis'' y Paso 3 ``Prueba''); Clase 1 --- Políticas y Control de
Acceso (\S OAuth 2.0 y \emph{scope}); Clase 3 --- Principios de Diseño y Vulnerabilidades
(\S Tipos de vulnerabilidades: relaciones de confianza, injection flaws; \S disponibilidad
y DoS como ataque de seguridad).

\paragraph{Qué necesitás saber.}
\begin{itemize}
  \item \textbf{Metodología de hipótesis de falla (memorizar):} Recolección de información
  $\to$ Planteo de hipótesis de falla $\to$ Prueba de la vulnerabilidad $\to$
  Generalización $\to$ Eliminación (opcional). Acá se pide directamente el paso 2
  (hipótesis) justificado con el paso 3 (cómo probarla): ``definir cómo probar'' prioriza
  \emph{analizar documentación o comportamiento} antes que intentar explotar, y la prueba
  debe ser \textbf{repetible y concluyente}.
  \item \textbf{Relaciones de confianza (Clase 3, tipos de vulnerabilidades):} ``no
  chequear autorización porque \emph{nadie va a adivinar la URL}'' es exactamente el
  patrón de la URL \texttt{/\{user\_id\}/public/calendar}: hay que verificar que el
  servidor efectivamente valide el flag de ``habilitado'' del usuario y no simplemente
  confíe en que el \texttt{user\_id} es difícil de adivinar.
  \item \textbf{OAuth y \emph{scope} (Clase 1):} el token de acceso debe limitarse al
  \emph{scope} mostrado en la pantalla de consentimiento (acá, lectura de reuniones);
  pedir de más (por ejemplo, permiso de escritura/administración total del calendario) es
  una falla de diseño de autorización.
  \item \textbf{Disponibilidad como propiedad de seguridad (Clase 3):} un ataque que tira
  o degrada el sistema (DoS) \emph{es} un ataque de seguridad aunque no comprometa
  confidencialidad ni integridad; un proceso de complejidad $O(n^3)$ sobre un dato que un
  atacante puede inflar (cantidad de eventos del mes) es candidato directo.
  \item \textbf{Injection flaws (Clase 3):} el profesor remarca que $\sim$90\% de las
  vulnerabilidades caen en esta familia; cualquier campo de entrada no controlado (acá, el
  ``mensaje'' que deja un desconocido en el calendario público) es un lugar concreto donde
  buscarla.
\end{itemize}

\paragraph{Notas y conexiones.}
Las cuatro hipótesis se ubican, cada una, en un componente concreto descripto en el
enunciado (la URL pública, el proceso $O(n^3)$, el flujo OAuth con Google/Apple, el campo
de mensaje libre), tal como pide la aclaración de no usar hipótesis genéricas. Todas están
tomadas de puntos ya vistos en el material (relaciones de confianza, disponibilidad/DoS,
scope de OAuth, injection flaws) y no se agrega nada fuera de eso.

\paragraph{Resolución.}

\textbf{Hipótesis 1 --- IDOR / relación de confianza en \texttt{/\{user\_id\}/public/calendar}.}
Es probable que el servidor, al recibir un \texttt{GET} a esa URL, no vuelva a validar que
el usuario efectivamente tenga habilitada la opción ``pública'' en su perfil, confiando en
que \texttt{user\_id} es difícil de adivinar o de enumerar (la relación de confianza
``nadie va a adivinar la URL'' que señala el material). Si los \texttt{user\_id} son
secuenciales o de algún modo predecibles, cualquiera podría ver/agendar en calendarios que
el dueño nunca hizo públicos.
\emph{Cómo probarla:} crear una cuenta de prueba con la página pública \textbf{deshabilitada},
y con otra cuenta (o sin sesión) acceder directamente a
\texttt{https://organizame.io/\{ese\_user\_id\}/public/calendar}; si igual devuelve datos
de agenda o permite agendar, queda demostrado. También probar con \texttt{user\_id}
consecutivos a uno propio habilitado, para ver si expone calendarios de terceros que no lo
sabían.

\textbf{Hipótesis 2 --- Denegación de servicio vía el proceso $O(n^3)$.} Es probable que el
proceso de análisis de complejidad $O(n^3)$ en la cantidad de eventos del mes no tenga un
límite de eventos por usuario/mes, por lo que un usuario (o un atacante que agenda a través
de la página pública de un tercero) podría inflar artificialmente la cantidad de eventos de
un mes para que ese proceso batch consuma recursos desproporcionados, degradando la
disponibilidad del servicio para otros usuarios.
\emph{Cómo probarla:} con una cuenta de prueba, generar mediante un script una gran
cantidad de eventos en un mismo mes (por API propia o a través del flujo de agendar
reunión desde una página pública), y medir el tiempo/CPU que insume la corrida del proceso
de análisis para ese mes, verificando si crece de forma cúbica y si llega a afectar al
resto del sistema (colas, timeouts, caída del servicio).

\textbf{Hipótesis 3 --- Scope de OAuth excesivo al integrar Google/Apple Calendar.} Es
probable que, al pedir ``acceso al calendario'' de la cuenta de Google/Apple, la app
solicite un \emph{scope} más amplio del necesario (por ejemplo, lectura \textbf{y}
escritura/administración del calendario, o acceso a otros datos de la cuenta), cuando para
leer reuniones agendadas alcanzaría con un scope de solo lectura.
\emph{Cómo probarla:} iniciar el flujo de registro con Google/Apple y revisar la pantalla
de consentimiento y/o decodificar el \emph{access token} obtenido, verificando qué scopes
concretos se solicitan y comparándolos con lo mínimo necesario (leer eventos); si el scope
permite modificar o borrar eventos de la cuenta real de Google/Apple del usuario sin que
eso sea necesario para la funcionalidad descripta, queda demostrado el exceso.

\textbf{Hipótesis 4 --- Injection en el campo de mensaje de la página pública.} Es probable
que el campo de texto donde un desconocido deja un mensaje al agendar una reunión desde
\texttt{/\{user\_id\}/public/calendar} no sanitice la entrada antes de guardarla o de
mostrarla en el calendario del dueño, permitiendo una inyección (por ejemplo, un payload de
script si el mensaje se renderiza en la vista web del calendario del dueño).
\emph{Cómo probarla:} desde la página pública, agendar una reunión de prueba dejando como
mensaje un payload típico de prueba (por ejemplo una etiqueta de script inocua, o una
comilla simple si el mensaje se guarda en una base de datos SQL), y luego revisar, desde la
cuenta del dueño, si el mensaje se ejecuta/se interpreta en vez de mostrarse como texto
plano, o si genera un error de base de datos que delate una inyección.

\begin{ojo}
No se incluyen hipótesis genéricas (``puede haber XSS en algún lado'', ``puede haber SQL
injection en algún formulario'') porque la consigna explícitamente lo prohíbe: cada
hipótesis de arriba está anclada a un componente y un flujo concreto descripto en el
enunciado de ``Organizame''.
\end{ojo}


\end{document}
